% Pseudo-Boolean encodings and proof logging for integer variables with
% several views.
%
% Build with: latexmk -pdf literal-encodings.tex
%
\documentclass[11pt,a4paper]{article}

\usepackage[a4paper,margin=2.6cm]{geometry}
\usepackage{amsmath,amssymb,amsthm}
\usepackage{mathtools}
\usepackage{booktabs}
\usepackage{enumitem}
\usepackage[T1]{fontenc}
\usepackage{lmodern}
\usepackage{microtype}
\usepackage{xcolor}
\usepackage[colorlinks,linkcolor=blue!50!black,citecolor=blue!50!black,urlcolor=blue!50!black]{hyperref}
\usepackage{cleveref}

% ---------------------------------------------------------------- theorem envs
% Numbering follows the thesis: results that are the thesis's results keep the
% thesis's numbers, and results that are new to this document are lettered.  See
% "About this document".
\theoremstyle{plain}
\newtheorem{theoremT}{Theorem}
\newtheorem{lemmaT}{Lemma}
\newtheorem{corollaryT}{Corollary}
\renewcommand{\thetheoremT}{3.\arabic{theoremT}}
\renewcommand{\thelemmaT}{3.\arabic{lemmaT}}
\renewcommand{\thecorollaryT}{3.\arabic{corollaryT}}

\newtheorem{theoremN}{Theorem}
\newtheorem{lemmaN}{Lemma}
\renewcommand{\thetheoremN}{3.\Alph{theoremN}}
\renewcommand{\thelemmaN}{3.\Alph{lemmaN}}

% Family versions of thesis theorems keep the thesis's number, with a prime.
\newtheorem{theoremP}{Theorem}
\renewcommand{\thetheoremP}{3.\arabic{theoremP}$'$}

\theoremstyle{definition}
\newtheorem{definitionT}{Definition}
\renewcommand{\thedefinitionT}{3.\arabic{definitionT}}
\newtheorem{exampleT}{Example}
\renewcommand{\theexampleT}{3.\arabic{exampleT}}
\newtheorem{procedureT}{Encoding Procedure}
\renewcommand{\theprocedureT}{3.\arabic{procedureT}}

\theoremstyle{remark}
\newtheorem*{remark}{Remark}
\newtheorem{remarkN}{Remark}
\renewcommand{\theremarkN}{3.\arabic{remarkN}}

\renewcommand{\theequation}{3.\arabic{equation}}

\crefname{theoremT}{Theorem}{Theorems}
\crefname{theoremP}{Theorem}{Theorems}
\crefname{theoremN}{Theorem}{Theorems}
\crefname{lemmaT}{Lemma}{Lemmas}
\crefname{lemmaN}{Lemma}{Lemmas}
\crefname{corollaryT}{Corollary}{Corollaries}
\crefname{definitionT}{Definition}{Definitions}
\crefname{procedureT}{Encoding Procedure}{Encoding Procedures}
\crefname{exampleT}{Example}{Examples}
\crefname{remarkN}{Remark}{Remarks}

% ---------------------------------------------------------------------- macros
% Every function, set and procedure name goes through \fname, and every
% constraint name through \globalcon, so that their fonts can be changed in one
% place rather than hunted for.
\newcommand{\fname}[1]{\ensuremath{\operatorname{#1}}}
\newcommand{\globalcon}[1]{\ensuremath{\mathrm{#1}}}
\newcommand{\BinEnc}{\fname{BinEnc}}
\newcommand{\Bnd}{\fname{Bnd}}
\newcommand{\Def}{\fname{Def}}
\newcommand{\Lin}{\fname{Lin}}
\newcommand{\Aux}{\fname{Aux}}
\newcommand{\Vars}{\fname{Vars}}
\newcommand{\ViewVars}{\fname{ViewVars}}
\newcommand{\ViewLink}{\fname{ViewLink}}
\newcommand{\Cells}{\fname{Cells}}
\newcommand{\Bdry}{\fname{Bd}}
\newcommand{\dm}{\fname{dom}}
\newcommand{\ImpLits}{\fname{ImpLits}}
\newcommand{\Standardise}{\fname{Standardise}}
\newcommand{\Linearise}{\fname{Linearise}}
\newcommand{\NeededViews}{\fname{NeededViews}}
\newcommand{\NeededAtomicVariables}{\fname{NeededAtomicVariables}}
\newcommand{\atge}[2]{#1_{\geq #2}}
\newcommand{\ateq}[2]{#1_{= #2}}
\newcommand{\atin}[3]{#1_{\in[#2,#3]}}
% Negated forms: the bar sits over the variable name only, never over the
% subscript, which otherwise stretches it unpleasantly.
\newcommand{\natge}[2]{\overline{#1}_{\geq #2}}
\newcommand{\nateq}[2]{\overline{#1}_{= #2}}
\newcommand{\natin}[3]{\overline{#1}_{\in[#2,#3]}}
% The atom of a set of values, for when the set matters more than its endpoints.
% A capital letter is always the set, never the atom over it.
\newcommand{\atset}[2]{#1_{\in #2}}
\newcommand{\natset}[2]{\overline{#1}_{\in #2}}
\newcommand{\inv}[1]{\textnormal{\textsc{Inv-#1}}}
\newcommand{\code}[1]{\texttt{#1}}
% Unnumbered sections are referred to by title, so that the numbered ones
% keep the thesis's numbering.
\newcommand{\nsec}[1]{\hyperref[#1]{``\nameref*{#1}''}}

% Invariant names, so a typo is a compile error rather than a silent
% disagreement between two paragraphs.
\newcommand{\invChain}{\inv{Chain}}
\newcommand{\invCut}{\inv{Cut}}
\newcommand{\invPart}{\inv{Part}}
\newcommand{\invCover}{\inv{Cover}}
\newcommand{\invCont}{\inv{Cont}}
\newcommand{\invView}{\inv{View}}
\newcommand{\invName}{\inv{Name}}
\newcommand{\invConflict}{\inv{Conflict}}
\newcommand{\invDomain}{\inv{Domain}}

\title{Representing Problems for Proofs, and a Proof Logging Framework,\\
  for Integer Variables with Several Views}
\author{Matthew McIlree \and Ciaran McCreesh}
\date{}

\begin{document}
\maketitle

\section*{About this document}

This is a revised and extended version of Sections~3.2 and~3.3 of the first
author's PhD thesis, \emph{Pseudo-Boolean Proof Logging for Constraint
Propagation Algorithms} (University of Glasgow, 2026), adapted and reused with
substantial additions.  It exists for three reasons: as the basis of a journal
treatment of pseudo-Boolean proof logging for constraint programming; as the
single authoritative account, for the Glasgow Constraint Solver's developers, of
what a literal is and what makes the literal layer work; and as a target for
adversarial review.  A subtle defect in this design is much cheaper to find now
than to diagnose later from an incomprehensible verifier rejection.

The thesis presents a canonical pseudo-Boolean representation of integer
variables using two's complement bit-vectors.  Additionally, there are two
families of ``atomic literals'': inequality atoms $\atge{x}{v}$ and equality
atoms $\ateq{x}{v}$.  The solver has since gained interval literals and views,
neither of which the thesis covers.

\begin{itemize}
  \item \textbf{Interval literals.}  There is now a third family, $\atin{x}{a}{b}$,
    asserting that $X$ takes a value in the closed interval $[a,b]$.  It exists so
    that a propagator can remove a run of values, and a reason can name a run of
    holes, at constant rather than linear cost in the width of the run.  A
    width-one interval literal \emph{is} the equality atom, not a separate
    Boolean.
  \item \textbf{Several views of one variable.}  A constraint may be
    posted over an affine \emph{view} $V = sX+c$ of a variable $X$, and one $X$
    may be wrapped by several views at once.  Each registered view is given its
    own \emph{view variable}, its own bit vector, and its own complete family of
    atomic literals, joined to $X$'s by a definitional link and by atom-level
    linking constraints.
\end{itemize}

Neither addition is conservative over the thesis's propagation results.
Inequality and equality atoms are complete under unit propagation given the
order chain alone; interval literals are not, and need two further families of
state-independent clause.  Facts move between two view variables of the same
variable only through linking constraints, and --- for interval literals --- the
direction that a proof actually needs is the one that is \emph{not} derivable
from the others.  The theorems below are therefore restated over a
\emph{view family} and over all three literal kinds, and reproved.

\paragraph{Numbering.}  Results that correspond to results of the thesis keep
the thesis's numbers, so that citations from the thesis and from our earlier
papers still land where they should: \cref{lem:composing,lem:contra-bounds,%
lem:chain,lem:reasons}, \cref{cor:chainrup}, and
\cref{thm:p1p2,thm:wipeout,thm:complete,thm:btrup,thm:maintain}.  Results that
are new to this document are lettered: \cref{lem:normalise,lem:contclosure,lem:cell,lem:covering,%
lem:transport}, and the family versions of \cref{thm:wipeout} and
\cref{thm:complete} are primed.  Invariants are named rather than numbered, with the
thesis's numbers given alongside, because an argument about \textsc{Inv-Cover}
is easier to follow than one about ``Inv5''.

\paragraph{Scope.}  Section~3.4 of the thesis, on justification procedures for
individual propagators, is unchanged by any of this and is not reproduced.  Two
points at which it meets the material here are flagged under
\nsec{sec:notprovided}.

\setcounter{section}{2}
\section{Proof Logging for Constraint Programming}

\subsection*{3.1\quad Constraint Programming Fundamentals}
\addtocounter{subsection}{1}

Unchanged; see the thesis.  We assume a constraint satisfaction problem
$P=\langle\mathcal V,\dm_0,\mathcal C\rangle$ with finite integer domains, and
a solver that alternates constraint propagation with backtracking search.

\subsection{Representing Problems for Proofs}
\label{sec:representing}

The philosophy of using pseudo-Boolean proof logging for constraint programming
solvers hinges on the idea that, for any CSP or COP, we can produce a
pseudo-Boolean encoding that can be reasonably understood as (and trusted to
be) the same problem.

\setcounter{exampleT}{2}
\begin{exampleT}[PB encoding of a CP problem]
\label{ex:pbenc}
Consider the CSP with variables $X,Y,Z$ all with domain $\{0,\dots,7\}$ and
constraints $X=Y$, $Y\neq Z$, $X+Z\geq 3$.  Intuitively, the PB formula
\begin{align}
  x_0+2x_1+4x_2-y_0-2y_1-4y_2 &\geq 0; \label{eq:ex1}\\
  -x_0-2x_1-4x_2+y_0+2y_1+4y_2 &\geq 0; \\
  u \Rightarrow y_0+2y_1+4y_2-z_0-2z_1-4z_2 &\geq 1;\\
  \bar u \Rightarrow -y_0-2y_1-4y_2+z_0+2z_1+4z_2 &\geq 1;\\
  x_0+2x_1+4x_2+z_0+2z_1+4z_2 &\geq 3 \label{eq:ex5}
\end{align}
represents the same problem, since we can view each of the PB variables $x_i$,
$y_i$ and $z_i$ as representing the value of the $i$-th bit in the binary
representation of the CP variables, and the ``flag'' variable $u$ ensures that
either $Y>Z$ or $Z>Y$.
\end{exampleT}

In this section we expand on the intuition of the above example and give a
method, following Gocht et al., for producing a PB formula that is obviously
equivalent to an input CSP, at least for some generalised definition of the
word ``obvious''.  For clarity, we will talk about transforming a ``CP
problem'' into a ``PB problem''.

\subsubsection{Defining Equivalence for Proofs}

The ultimate goal of encoding for proof logging purposes is to trust that the
process preserves the integrity of the proof with respect to the original
problem, and potentially to verify this process formally in a theorem prover.
So at one level, the only guarantees we really need are:

\begin{description}[font=\normalfont\bfseries,leftmargin=2.4em]
  \item[P1.] The PB problem is satisfiable if and only if the CP problem is
    satisfiable.
  \item[P2.] If there is an objective variable, the optimal value of the CP
    objective is equal to the optimal value of the PB objective.
\end{description}

However, to make it easier to obtain these guarantees, we should employ very
restricted transformations with minimal reformulation, much less than one might
use if encoding in order to use a PB or SAT system to solve the problem.  To put
it simply, we want our encodings to be as ``dumb'' as possible because we want to
be extremely certain the two properties hold.  So in practice, our working
definition of equivalence is: ``the PB problem is the result of applying
\cref{proc:encoding} to the CP problem''.  We will define this procedure shortly.

The one substantive change to this section from the thesis is that the encoding
now admits two further \emph{shared} classes of introduced variable, beyond the
atomic constraint variables: \emph{view variables} for affine views,
and a third kind of atomic variable for intervals.  Both are still uniquely
determined, semantically, by an assignment to the original variables: their
definition constraints and their links leave no freedom.  That is what keeps P1
and P2 easy to argue.

\subsubsection{A Transformation Procedure}
\label{sec:transformation}

The basic idea is to first turn the CP problem into an equivalent problem with
only reified linear constraints (an ``intermediate ILP'').  This still contains
the original non-Boolean variables of the problem, so it is easy to argue
equivalence by showing that the effect of the linear constraints is to restrict
the variables' values in exactly the same way as the original CP constraints.
We have the freedom to introduce auxiliary variables at this stage to assist in
the linearisation (such as the $u$ flag in \cref{ex:pbenc}).

Two classes of introduced variable are \emph{shared} between linearisations
rather than private to one, and each is uniquely determined by an assignment to
the original variables:

\begin{itemize}
  \item \textbf{View variables} name affine views $sX+c$ of an
    original variable.  A constraint posted over a view of $X$ is linearised
    over the corresponding view variable.
  \item \textbf{Atomic constraint variables} (or simply \emph{atomic
    variables}) are $0$--$1$ variables denoting basic relationships between a
    variable and a value or an interval of values.
\end{itemize}

Both come with \emph{definition constraints}, which are fully reified on fresh
variables.  That makes them redundant in one specific sense: the values of the
introduced variables are uniquely determined for any solution, so the solution
set projected onto the original variables is unchanged, and in particular
adding them cannot change satisfiability or the objective value.  It does not
make them dispensable, since they are what the proof reasons with; nor does it
mean that solutions can be counted without regard to them, which is what the
preserved set of \cref{sec:enumeration} is for.

All other variables introduced as part of a linearisation should not appear in
other linearisations.  The soundness of concatenating linearisations rests on
exactly that division; see \cref{lem:composing}.

We assume we start with a general constraint satisfaction problem, with
variable set $\mathcal V$, finite domains $\dm_0(X_i)$ for $X_i\in\mathcal V$,
and constraints $\mathcal C$.  All domains can be assumed to be ranges of
integers without holes, since any missing values can be modelled with
additional constraints if necessary.  We also assume that the constraints are a
flat list, without constructs such as loops that might be found in a
higher-level modelling language.

\begin{procedureT}
\label{proc:encoding}
For a CSP with variables $\mathcal V$, initial domain state (without holes)
$\dm_0$, and constraints $\mathcal C$, we can produce a PB formula as follows.
\begin{enumerate}[label=\arabic*:,leftmargin=2.6em,itemsep=1pt]
  \item $\mathcal W \leftarrow \bigcup_{X\in\mathcal V}\ViewVars(X)$
  \item $\ViewLink(\mathcal W\setminus\mathcal V) \leftarrow
        \bigcup_{W\in\mathcal W\setminus\mathcal V}\ViewLink(W)$
  \item $\Bnd(\mathcal W) \leftarrow \bigcup_{W\in\mathcal W}\Bnd(W,\dm_0)$
  \item $\mathcal A \leftarrow \bigcup_{C\in\mathcal C}
        \NeededAtomicVariables(C,\dm_0)$
  \item $\Def(\mathcal A) \leftarrow \bigcup_{y\in\mathcal A}\Def(y)$
  \item $\Lin(\mathcal C) \leftarrow \bigcup_{C\in\mathcal C}\Linearise(C,\dm_0)$
  \item $F \leftarrow \BinEnc\bigl(\Bnd(\mathcal W)\cup\ViewLink(\mathcal W\setminus\mathcal V)
        \cup\Def(\mathcal A)\cup\Lin(\mathcal C)\bigr)$
  \item $F \leftarrow \Standardise(F)$
\end{enumerate}
\end{procedureT}

Steps~1 and~2 are new; step~5 now has a third case.  The remainder of this
section gives more detail on each stage along with soundness arguments, and
introduces notation used throughout.

\paragraph{Steps 1 and 2: views and their links.}

\begin{definitionT}[View and view variable]
\label{def:view}
A \emph{view} of a CP variable $X$ is an affine expression $sX+c$ with
$s\in\{+1,-1\}$ and $c\in\mathbb Z$, and it is \emph{proper} when it is not $X$
itself.  The encoding gives each view its own \emph{view variable}: an integer
variable $W$ with a sign $s_W\in\{+1,-1\}$ and an offset $c_W\in\mathbb Z$,
subject to $W = s_W X + c_W$.  Write
\[
  \varphi_W(k) = s_W k + c_W, \qquad
  \psi_W(v) = \varphi_W^{-1}(v) = s_W(v - c_W),
\]
so that $\varphi_W$ is an order-preserving bijection $\mathbb Z\to\mathbb Z$
when $s_W=+1$ and an order-reversing one when $s_W=-1$.  Every CP variable is
the view variable of the trivial view of itself, with $\varphi_X$ the identity
function.  A view variable $W\neq X$ is \emph{proper}, and $X$ is its
\emph{base}.

Write $\NeededViews(C)$ for the views occurring in a constraint $C$, and
$\ViewVars(X)$ --- the \emph{view family} of $X$ --- for $X$ together with the
view variable of each proper view of $X$ in
$\bigcup_{C\in\mathcal C}\NeededViews(C)$.  Then
$\mathcal W = \bigcup_{X\in\mathcal V}\ViewVars(X)$, so $\ViewVars$ is a
partition of $\mathcal W$ by base, and $\mathcal W\setminus\mathcal V$ is the
set of proper view variables.
\end{definitionT}

For each proper $W\in\ViewVars(X)$, $\ViewLink(W)$ is the single linear equality
\begin{equation}
  \ViewLink(W) \;:=\; W - s_W X = c_W ,
  \label{eq:link}
\end{equation}
which \Standardise{} will split into a $\geq$ and a $\leq$ row.  It is stated
only for proper view variables: a base variable has nothing to link to.

View variables arise because constraints are posted over \emph{operands}, and an
operand may be a view.  A modelling layer that writes
$\globalcon{AllDifferent}(X_1, X_2+1, X_3+2)$ is naming three operands, two of
which are proper views of $X_2$ and $X_3$.  Two different constraints may name
the same view, and one variable may be named through several different views at
once, which is precisely why view variables have to be shared rather than
private to a linearisation.

A view therefore gets a variable of its own, at the cost of one linear equality.
\Cref{sec:rationale} says why, and what the alternative --- substituting
$sX+c$ wherever the view occurs --- would cost.

\paragraph{Step 3: domain bound constraints.}
For a view variable $W$ of $X$ and domain state $\dm_0$, let $\Bnd(W,\dm_0)$
be the two linear constraints
\begin{equation}
  W \geq \min \varphi_W(\dm_0(X)) \qquad\text{and}\qquad
  W \leq \max \varphi_W(\dm_0(X)) ,
  \label{eq:bnd}
\end{equation}
that is, $W\geq \varphi_W(l)$ and $W \leq \varphi_W(u)$ when $s_W=+1$, and the
two swapped when $s_W=-1$, writing $l$ and $u$ for the bounds of $\dm_0(X)$.
When it is clear from context we omit $\dm_0$ and write $\Bnd(W)$.  We call
$[\,\min\varphi_W(\dm_0(X)),\ \max\varphi_W(\dm_0(X))\,]$ the \emph{definition
range} of $W$ and denote it $[l_W,u_W]$.

Note that $\Bnd(W)$ for a proper $W$ is implied by $\Bnd(X)$ together with
$\ViewLink(W)$ (it is their sum, up to sign), so asserting it changes nothing.  It
is asserted anyway, because unit propagation cannot add two constraints
together, and several later arguments need $W$'s own bounds to be available as
units on $W$'s own atoms.

Clearly, if we construct a CSP with variables $\mathcal W$, constraints
$\mathcal C \cup \Bnd(\mathcal W)\cup\ViewLink(\mathcal W\setminus\mathcal V)$, and unrestricted
integer domains, then its solution set has the same cardinality as our original
CSP's, and is identical to it under projection onto $\mathcal V$: every proper
view variable is determined by its base through \cref{eq:link}, so the
projection is a bijection and not merely a surjection.

\paragraph{Steps 4 and 5: atomic constraint variables.}
These are fresh $0$--$1$ variables that can be shared between linearisations
and have a special meaning.  For each supported constraint type $C$ we need to
know which atomic variables it requires, so $\NeededAtomicVariables$ is
defined on a per-constraint basis; usually this is immediately clear from the
definition of $\Linearise(C,\dm_0)$.

Each atomic variable $y$ requires a fully reified linear definition constraint
$\Def(y)$ of the form $y\Leftrightarrow C$.  We allow three types of atomic
variable definition, each associated with a \emph{view variable}
$W\in\mathcal W$ --- not merely with a CP variable --- and with one or two
integers.  For clarity we always use subscripted names suggestive of the
definition, writing $w$ for the PB name of the view variable $W$.

\begin{enumerate}[leftmargin=2.2em]
  \item A \textbf{bound variable} $\atge{w}{v}$ represents whether $W$ must be
    at least $v$:
    \begin{equation}
      \Def(\atge{w}{v}) := \atge{w}{v} \Leftrightarrow W \geq v .
      \label{eq:defge}
    \end{equation}
  \item An \textbf{equality variable} $\ateq{w}{v}$ represents whether $W$ must
    equal $v$:
    \begin{equation}
      \Def(\ateq{w}{v}) := \ateq{w}{v} \Leftrightarrow
        \atge{w}{v} + \natge{w}{v+1} \geq 2 .
      \label{eq:defeq}
    \end{equation}
  \item An \textbf{interval variable} $\atin{w}{a}{b}$, for $a<b$, represents
    whether $W$ must take a value in the closed interval $[a,b]$:
    \begin{equation}
      \Def(\atin{w}{a}{b}) := \atin{w}{a}{b} \Leftrightarrow
        \atge{w}{a} + \natge{w}{b+1} \geq 2 .
      \label{eq:defin}
    \end{equation}
\end{enumerate}

Naturally, if an equality or interval variable is needed then so are the
corresponding bound variables.

An interval variable is a \emph{wide equality atom}: \cref{eq:defeq,eq:defin}
are the same definition, cut at two values instead of at $v$ and $v+1$.  This
is not a coincidence to be tidied away later, and we record the consequence
immediately, because it is the single most expensive mistake available here.

\begin{definitionT}[Width-one identification]
\label{def:widthone}
For $a=b$, the interval variable $\atin{w}{a}{b}$ \emph{is} the equality
variable $\ateq{w}{a}$: the same PB variable, not a second one with the same
definition.  For $a>b$, the interval condition is the constant false.
\end{definitionT}

A separate width-one interval flag would be an unlinked doppelg\"anger of the
equality atom: two Booleans with identical definitions, and nothing in the
formula connecting them, so that a fact recorded in one is invisible through
the other.  It cannot be repaired by adding the equivalence after the fact ---
see \nsec{sec:loadbearing} below, where this is witness~W1.

Because these are fully reified constraints on fresh variables, they are
entirely redundant: the values of the atomic variables are uniquely determined
for any solution.  This means the solution set projected onto the original
variables remains exactly the same, and in particular adding them cannot
possibly change satisfiability or the objective value.  The same is true of
view variables, by \cref{eq:link}.

Recalling $0$--$1$ literals, we can represent $W$ being at most $v$, not equal
to $v$, or outside $[a,b]$ by negating the corresponding variables.
Collectively we refer to literals of the form $\atge{w}{v}$,
$\natge{w}{v}$, $\ateq{w}{v}$, $\nateq{w}{v}$,
$\atin{w}{a}{b}$ and $\natin{w}{a}{b}$ as \emph{atomic (constraint)
literals}.

\paragraph{Step 6: linearise constraints.}
For each constraint $C$ we require a function $\Linearise(C,\dm_0)$ turning it
into a set of reified linear constraints $\Lin(C)$.  This can use
view variables, atomic literals, and further auxiliary variables.
Write $\Vars(\Lin(C))$ for the set of variables appearing in the
linearisation, and let $\Aux(C) = \Vars(\Lin(C))\setminus(\mathcal W\cup
\mathcal A)$ be the auxiliary variables introduced.  We require:

\begin{description}[font=\normalfont\bfseries,leftmargin=2.4em]
  \item[L1.] Each constraint in any $\Lin(C)$ fits one of two general forms:
    either
    \begin{equation}
      \sum_{j=0}^{n-1} a_j\cdot V_j \bowtie b
      \qquad\text{or}\qquad
      r_1\wedge\cdots\wedge r_k \rightleftharpoons \sum_{j=0}^{n-1} a_j\cdot V_j \bowtie b
      \label{eq:forms}
    \end{equation}
    where $b$ and each $a_j$ are integers; each $r_k$ is a $0$--$1$ literal;
    each $V_j$ is either a $0$--$1$ or an integer variable; $\rightleftharpoons$
    is one of $\{\Leftarrow,\Rightarrow,\Leftrightarrow\}$; and $\bowtie$ is one
    of $\{\geq,\leq,=\}$.
  \item[L2.] Any non-$0$--$1$ auxiliary integer variable $V_j\in\Aux(C)$ must be
    bounded with corresponding bound constraints $\Bnd(V_j)$, considered part of
    the linearisation.
  \item[L3.] Auxiliary variables must not be shared between linearisations:
    $\Aux(C_1)\cap\Aux(C_2)=\emptyset$ for all $C_1\neq C_2$.  View variable
    variables and atomic variables are \emph{not} auxiliary variables and may be
    shared freely.
  \item[L4.] Most importantly, $\Lin(C)\cup\Def(\mathcal A)\cup\ViewLink(\mathcal W\setminus\mathcal V)$
    must enforce the same restriction as $C$ on the variables in $\mathcal V$.
    That is, a valid assignment $\sigma:\mathcal V\to\mathbb Z$ satisfies $C$ if
    and only if it can be extended to a valid assignment
    $\sigma^+:\Vars(\Lin(C)\cup\Def(\mathcal A)\cup\ViewLink(\mathcal W\setminus\mathcal V))\to\mathbb Z$
    satisfying $\Lin(C)\cup\Def(\mathcal A)\cup\ViewLink(\mathcal W\setminus\mathcal V)$.
\end{description}

L3 is where view variables differ from the thesis's presentation:
they are shared, and they are shared for the same reason atomic variables are.
The proof of \cref{lem:composing} shows that this is safe precisely because
\cref{eq:link}, like \cref{eq:defge,eq:defeq,eq:defin}, determines the shared
variable's value uniquely from the assignment to $\mathcal V$.

\begin{lemmaT}[Composing linearisations]
\label{lem:composing}
Let $P$ be a CSP with variables $\mathcal V$ and constraints $\mathcal C$,
where all domains are integer ranges.  After executing steps~1 to~6 of
\cref{proc:encoding} let $P'$ be the CSP with constraints
$\mathcal C' := \Lin(\mathcal C)\cup\Bnd(\mathcal W)\cup\ViewLink(\mathcal W\setminus\mathcal V)
\cup\Def(\mathcal A)$, whose variables are precisely those occurring in
$\mathcal C'$, and whose domain for every non-$0$--$1$ variable is $\mathbb Z$.
Then:
\begin{enumerate}[label=\arabic*.,itemsep=0pt]
  \item every solution $\sigma$ of $P$ can be extended to a solution $\sigma'$
    of $P'$;
  \item every solution $\sigma'$ of $P'$ can be restricted to a solution
    $\sigma$ of $P$.
\end{enumerate}
\end{lemmaT}

\begin{proof}
Call a variable of $P'$ \emph{determined} if its value is a function of the
restriction of an assignment to $\mathcal V$ alone.  Every variable in
$\mathcal W\cup\mathcal A$ is determined: a proper view variable by
\cref{eq:link}, and an atomic variable by its definition constraint, which is a
fully reified constraint over $\mathcal W$ (for $\atge{w}{v}$) or over
previously determined atomic variables (for $\ateq{w}{v}$ and $\atin{w}{a}{b}$,
whose definitions mention only bound variables of the same view variable).
The determination is well founded because the dependency order is
$\mathcal V \to \mathcal W \to \{\text{bound}\} \to
\{\text{equality},\text{interval}\}$, with no cycles.

\emph{Part 1.}  Let $\sigma$ be a solution to $P$ and $C_1,\dots,C_n$ the
constraints of $P$.  By validity, $\sigma$ satisfies $\Bnd(\mathcal V)$, and
extending it by the determined values satisfies $\Bnd(\mathcal
W)\cup\ViewLink(\mathcal W\setminus\mathcal V)$.  By L4, $\sigma$ can be extended to a solution
$\sigma_1$ of $\Bnd(\mathcal W)\cup\ViewLink(\mathcal W\setminus\mathcal V)\cup\Def(\mathcal
A)\cup\Lin(C_1)$.

Now assume we have extended $\sigma$ to a solution $\sigma_{k-1}$ of
$\Bnd(\mathcal W)\cup\ViewLink(\mathcal W\setminus\mathcal V)\cup\Def(\mathcal A)\cup
\bigcup_{i<k}\Lin(C_i)$.  By L4 again, $\sigma$ can be extended to a solution
$\sigma'_k$ of $\Bnd(\mathcal W)\cup\ViewLink(\mathcal W\setminus\mathcal V)\cup\Def(\mathcal
A)\cup\Lin(C_k)$.  Define $\sigma_k$ by
\begin{align}
  \sigma_k(V) &= \sigma(V) && \text{if } V\in\mathcal V; \\
  \sigma_k(V) &= \sigma_{k-1}(V) && \text{if } V\in\mathcal W\cup\mathcal A; \\
  \sigma_k(V) &= \sigma_{k-1}(V) && \text{if otherwise }
     V\in\Vars\bigl(\bigcup_{i<k}\Lin(C_i)\bigr); \\
  \sigma_k(V) &= \sigma'_k(V) && \text{if otherwise } V\in\Vars(\Lin(C_k)).
\end{align}
This is an extension of $\sigma$, and it satisfies both $\Lin(C_k)$ and
$\Bnd(\mathcal W)\cup\ViewLink(\mathcal W\setminus\mathcal V)\cup\Def(\mathcal A)\cup
\bigcup_{i<k}\Lin(C_i)$, because
$\sigma_k(V)=\sigma_{k-1}(V)=\sigma'_k(V)$ for every $V\in\mathcal
V\cup\mathcal W\cup\mathcal A$ --- each such $V$ being determined by $\sigma$,
which all three assignments extend --- while
$\Vars(\bigcup_{i<k}\Lin(C_i))\cap\Vars(\Lin(C_k))\subseteq
\mathcal V\cup\mathcal W\cup\mathcal A$ by L3, so $\sigma_k$ is an extension of
both $\sigma_{k-1}$ and $\sigma'_k$.  The claim follows by induction.

\emph{Part 2.}  Let $\sigma'$ be a solution to $P'$.  There is only one
possible restriction of $\sigma'$ to a valid assignment to $\mathcal V$; and by
L4, since $\sigma'$ satisfies $\Lin(C_i)\cup\Def(\mathcal A)\cup\ViewLink(\mathcal
W\setminus\mathcal V)$ for each $i$, that restriction must satisfy each $C_i$ and hence $P$.
\end{proof}

We refer to the result of steps~1 to~6 as the \emph{intermediate ILP}.  In the
linearisations presented here, L1 to L4 are verifiable by inspection.

\begin{exampleT}
An intermediate ILP for the CSP of \cref{ex:pbenc}, where the third constraint
is posted over the view $Z' = -Z + 7$ as $X - Z' \geq -4$, could be
\begin{align}
  &X,Y,Z,Z',u\in\mathbb Z ;\\
  &0\leq X\leq 7;\quad 0\leq Y\leq 7;\quad 0\leq Z\leq 7;\quad
    0\leq Z'\leq 7;\quad 0\leq u\leq 1;\\
  &Z' + Z = 7; \label{eq:exlink}\\
  &X - Y = 0;\\
  &u \Rightarrow X - Y > 0; \qquad \bar u \Rightarrow Y - X > 0;\\
  &X - Z' \geq -4 .
\end{align}
Line \cref{eq:exlink} is $\ViewLink(Z')$, and $\Bnd(Z')$ is the image of
$\Bnd(Z)$ under $\varphi_{Z'}(k) = -k+7$.
\end{exampleT}

\paragraph{Step 7: binary encoding of integer variables.}
$\BinEnc$ turns the intermediate ILP into a $0$--$1$ ILP by replacing each
occurrence of a non-Boolean variable with a binary exponential sum.  This may
be a poor encoding from a solving perspective, but it maintains the form of the
ILP, letting us ``pretend'' we have integer variables while working honestly
with native PB constraints.  Every variable in $\mathcal W$ --- proper
view variables included --- is replaced, and each gets its \emph{own} bit
vector: $\BinEnc(W)$ and $\BinEnc(X)$ share no PB variables, and are related
only by the encoded form of $\ViewLink(W)$.

Specifically, for any non-$0$--$1$ variable $W$ with definition range
$[l_W,u_W]$: if $l_W\geq 0$, let $h$ be least with $2^h\geq u_W+1$ and define
$\BinEnc(W) := \sum_{i=0}^{h-1}2^i w_i$; otherwise let $h$ be least with
$2^{h-1}\geq \max\{|u_W|+1,|l_W|\}$ and define
$\BinEnc(W) := -2^{h-1}w_{h-1}+\sum_{i=0}^{h-2}2^i w_i$, the two's-complement
representation.  We overload the notation and write $\BinEnc(C)$ and
$\BinEnc(F)$ for the results of replacing every eligible variable in a
constraint or a set of constraints.

The soundness of this replacement follows directly from the validity of binary
and two's-complement encodings: there is a one-to-one correspondence between
the solutions of the $0$--$1$ ILP and those of the intermediate ILP, given by
the encoding.

\begin{remark}
Two width details matter in practice and neither affects the argument.  First,
the bit vector of a proper view variable is sized from $[l_W,u_W]$, and can be
\emph{wider} than that of its base: a view of a variable with range $[1,8]$ has
range $[1,8]$ too but a bit vector spanning $[0,15]$, and a view $X+10^6$ has a
much wider one.  Any coefficient bound or big-$M$ computed for a line mentioning
$W$ must be measured against $\BinEnc(W)$, not against the value range of $X$.
Second, the implementation's width rule diverges from Step~7 in two places, both
deliberate and both for conformance with the verified encoder
\code{cake\_pb\_cp}.  It allocates no magnitude bit that cannot reach a value of
the domain, so a variable with range $[0,0]$ has the empty sum, identically zero,
where the thesis's ``least strictly positive $h$'' gives it one bit.  And it
sizes from the \emph{signed} upper bound, $\max\{|l_W|,\,u_W+1\}$ rather than
$\max\{|u_W|+1,\,|l_W|\}$, which differs exactly for a wholly-negative singleton
at a power of two: $\{-2\}$ gets two bits rather than three, $\{-8\}$ four rather
than five.  Both encodings cover the definition range, so \cref{thm:p1p2} is
indifferent; but a view $-X$ of a variable fixed at a power of two meets the
second one.
\end{remark}

\paragraph{Step 8: standardise constraints.}
Finally we ensure all remaining constraints are in canonical form: convert
reifications to native PB constraints, replace $=$ constraints with a pair of
$\geq$ and $\leq$ constraints, and multiply through by $-1$ and/or add $1$ as
needed so that all inequalities are in $\geq$ form.  Such transformations do not
affect the solution set.

\begin{theoremT}
\label{thm:p1p2}
Let $P$ be a CSP with variables $\mathcal V$ and constraints $\mathcal C$, where
all domains are integer ranges.  Let $F$ be the result of executing
\cref{proc:encoding} on $P$.  Additionally, if $P$ has an associated objective
variable $X_o\in\mathcal V$ to minimise or maximise, let $\BinEnc(X_o)$ be the
PB objective function.  Then P1 and P2 hold.
\end{theoremT}

\begin{proof}
Let $P'$ be the intermediate ILP produced after steps~1 to~6.
\cref{lem:composing} gives a bijection between the solutions of $P$ and the set
of assignments to $\mathcal V$ (projected from full assignments) that satisfy
$P'$.  Let $F^*$ be the $0$--$1$ linear CSP obtained after step~7.  Every
solution to $P'$ corresponds to a solution of $F^*$, by assigning to each
non-$0$--$1$ variable $W$ the unique $h$-bit vector representing its value ---
note that this is well defined for a proper view variable exactly as for a
base one, since $\ViewLink(W)$ and $\Bnd(W)$ force $W$'s value into $[l_W,u_W]$,
which its bit width covers.  Conversely, every solution to $F^*$ corresponds to
a solution of $P'$ by decoding each bit vector.  The final formula $F$ obtained
after step~8 has exactly the same solutions as $F^*$.

So there is a bijection $\phi$ between the solutions of $P$ and the solutions of
$F$ projected onto the bit variables appearing in $\BinEnc(X)$ for
$X\in\mathcal V$.  In particular $F$ is satisfiable iff $P'$ is satisfiable iff
$P$ is satisfiable, so P1 holds; and for all $X\in\mathcal V$ and all solutions
$\sigma$ to $P$ we have $\sigma(X)=\BinEnc(X)\!\restriction_{\phi(\sigma)}$, so
P2 holds.
\end{proof}

\begin{remark}
The projection in \cref{thm:p1p2} is onto the bits of the \emph{base} variables
only.  The bits of proper view variables, and all atomic variables, are
determined but not projected onto.  This matters for enumeration; see
\cref{sec:enumeration}.
\end{remark}

\subsubsection{PB Encodings for Fundamental Constraints}

To instantiate the encoding method for concrete CSPs we exhibit ``sufficiently
obvious'' linearisations for all supported non-linear constraints.  These are
unchanged from the thesis, and are not reproduced here; the encodings for
$\globalcon{NotEquals}$, $\globalcon{Abs}$, $\globalcon{Element}$,
$\globalcon{ArrayMin}$, $\globalcon{ArrayMax}$, $\globalcon{Count}$,
$\globalcon{NValue}$, $\globalcon{AllDifferent}$ and $\globalcon{Table}$
satisfy L1 to L4 by inspection.

Two remarks are needed in the extended setting.

First, a constraint's operands may be views rather than base variables, and its
linearisation is written over their view variables.  Thus
$\globalcon{AllDifferent}(W_1,\dots,W_n)$ linearises to
$u_{ij}\Rightarrow W_i - W_j > 0$ and $\bar u_{ij} \Rightarrow W_j - W_i > 0$
over exactly the operands it was given, with no reference to their bases.  This
is what ``the encoding is generic in the operand'' means, and it is the property
that \cref{sec:framework} preserves at the level of atoms.

Second, no linearisation mentions interval variables.  Interval literals do not
appear in the model at all: they are introduced inside the proof, by redundance,
as the proof needs them.  The one place where the question arises --- whether
holes in an initial domain should be stated in the model as
$\natin{x}{a}{b}$ rather than derived at the root of the proof --- is
an open decision recorded in \cref{sec:open}.

\subsection{A Proof Logging Framework}
\label{sec:framework}

We begin from the assumption that we have used an encoding method following
\cref{proc:encoding}.  So before logging anything in the proof, we have a
correct PB representation of an input CP problem in terms of binary-encoded
view variables of the CP variables, atomic literals defined for certain view
variables and values, and some other auxiliary variables.

Based on this encoding method, we formulate a proof logging framework based on
maintaining invariants relating the solver's explored domain states to the
verifier's PB constraint database.  Certain properties of atomic literals are
crucial for this setup, so we demonstrate these first.

Throughout, ``UP'' means plain unit propagation from the empty assignment, with
no case splitting: the propagation a RUP check performs after asserting the
negation of the line being checked.  We write $F\!\restriction_R$ for the
formula $F$ restricted by the literals of $R$, and say that a literal $\ell$
\emph{propagates under $R$} if $\ell\in R$ or unit propagation of
$F\!\restriction_R$ sets $\ell$.

\subsubsection{Properties of Atomic Literals}
\label{sec:atomprops}

\paragraph{Definitions in the proof.}
We can always freely introduce in the proof further atomic variables with
definitions corresponding to \cref{eq:defge,eq:defeq,eq:defin}, by
redundance-based strengthening.  After binary encoding and standardisation
these amount, for a view variable $W$ of $X$, to the PB constraints
\begin{align}
  \Def^{\Rightarrow}(\atge{w}{v}) &:= \atge{w}{v} \Rightarrow \BinEnc(W)\geq v;
    \label{eq:dge1}\\
  \Def^{\Leftarrow}(\atge{w}{v}) &:= \natge{w}{v} \Rightarrow
    -\BinEnc(W)\geq -v+1; \label{eq:dge2}\\
  \Def^{\Rightarrow}(\ateq{w}{v}) &:= \ateq{w}{v} \Rightarrow
    \atge{w}{v} + \natge{w}{v+1} \geq 2; \label{eq:deq1}\\
  \Def^{\Leftarrow}(\ateq{w}{v}) &:= \nateq{w}{v} \Rightarrow
    \natge{w}{v} + \atge{w}{v+1} \geq 1; \label{eq:deq2}\\
  \Def^{\Rightarrow}(\atin{w}{a}{b}) &:= \atin{w}{a}{b} \Rightarrow
    \atge{w}{a} + \natge{w}{b+1} \geq 2; \label{eq:din1}\\
  \Def^{\Leftarrow}(\atin{w}{a}{b}) &:= \natin{w}{a}{b} \Rightarrow
    \natge{w}{a} + \atge{w}{b+1} \geq 1. \label{eq:din2}
\end{align}

We assume in what follows that any PB literal appearing in a derivation with
these suggestive names is accompanied by the necessary definition constraints,
either in the input formula or derived by redundance at an earlier point.  When
we say an atomic literal $\ell$ is \emph{defined} at a particular point in the
proof, we mean that both directions of the reified definition for its
underlying atomic variable have appeared in the accumulated formula and have
not been deleted.

Call the equality and interval atoms the \emph{two-cut} atoms, since each is
defined by \cref{eq:defeq,eq:defin} from a pair of bound variables, and their
literals the \emph{two-cut} literals.  From here on, when a two-cut atom is
discussed through the set of values it names we write $\atset{w}{S}$ for the
atom of $W$ that holds exactly when $W$ takes a value in the interval $S$, so
that $\atset{w}{[a,b]}$ is $\atin{w}{a}{b}$ for $a<b$ and $\atset{w}{\{v\}}$ is
$\ateq{w}{v}$ by \cref{def:widthone}.  In everything below, $E$, $C$, $G$ and
$S$ denote sets of integers and never the atoms over them: $E\subseteq E'$,
$E\cap[l_W,u_W]$ and $\bigcup_i C_i$ mean what they say about integers, while
$\atset{w}{E}$ and $\natset{w}{E}$ are the two literals of the atom naming $E$.

Note the asymmetry between
$\Def^{\Rightarrow}$ and $\Def^{\Leftarrow}$ for them,
\cref{eq:deq1,eq:deq2,eq:din1,eq:din2}.  The forward direction is
a \emph{conjunction}: asserting $\ateq{w}{v}$ or $\atin{w}{a}{b}$ makes both
cuts units.  The reverse direction is a \emph{disjunction}: asserting
$\nateq{w}{v}$ or $\natin{w}{a}{b}$ makes neither cut a
unit; \cref{eq:deq2,eq:din2} become binary clauses that fire only when one of
the two cuts is already known.  Almost everything in this section that is not
already in the thesis exists because of that asymmetry.

\begin{remark}
The solver optionally uses a compact encoding in which an equality atom at a
definition bound is defined by its one non-trivial cut only
($\ateq{w}{l_W}\Leftrightarrow\natge{w}{l_W+1}$, and dually at
$u_W$).  This is off by default, because the eager form matches
\code{cake\_pb\_cp}'s.  Under that variant \invCut{} \emph{fails} at the
definition bounds: the dropped cut is not a pinned atom but no atom at all ---
for $X\in[3,9]$ the equality atom at $3$ is defined and $\atge{x}{3}$ is never
created.  Everything below is therefore stated for the default encoding.  The
compact variant should be sound, the dropped cut being a constant of the
encoding rather than a fact, but making the results unconditional would mean
rereading every appeal to \invCut{} at a definition bound with a constant in
place of the atom, and we have not done that.
\end{remark}

\begin{lemmaT}[Propagation of contradictory bound literals]
\label{lem:contra-bounds}
Suppose two bound variables $\atge{w}{v}$ and $\atge{w}{v'}$ are defined for
some view variable $W$ and values $v\geq v'$.  Then if $\atge{w}{v}$ and
$\natge{w}{v'}$ propagate, further unit propagation will result in a
conflict.
\end{lemmaT}

\begin{proof}
Directly from Theorem~2.7 of the thesis, since even if $W$ requires a
two's-complement encoding, the definition constraints
$\Def^{\Rightarrow}(\atge{w}{v})$ and $\Def^{\Leftarrow}(\atge{w}{v'})$ are in
the required form after substituting $\atge{w}{v}\mapsto 1$,
$\atge{w}{v'}\mapsto 0$.  Note that both constraints are over the \emph{same}
bit vector $\BinEnc(W)$; this is why a view variable is given its own bit
vector rather than being unfolded into its base's.
\end{proof}

\begin{corollaryT}[RUP properties of atomic literals]
\label{cor:chainrup}
For two bound variables $\atge{w}{v}$ and $\atge{w}{v'}$ defined on the same
view variable $W$ with $v\geq v'$, the constraint
\begin{equation}
  \atge{w}{v} \Rightarrow \atge{w}{v'} \geq 1
  \label{eq:chainlink}
\end{equation}
always follows by RUP.
\end{corollaryT}

\paragraph{The invariants of the literal layer.}
We now state the invariants that the proof must maintain for the literal layer
to behave.  They are all \emph{state independent}: each is a tautological
consequence of the definitions, emitted once, at the top proof level, with
nothing to undo on backtracking.  We name them rather than number them; the
first, and the two in \cref{sec:unsat}, are Inv1, Inv2 and Inv3 of the thesis.

Fix a base variable $X$, and let $F_t$ be the constraints accumulated at a point
$t$ in the proof.

\begin{description}[font=\normalfont\bfseries,leftmargin=3.2em,style=nextline]

\item[\invChain{} (thesis Inv1): the order chain.]
  For any two bound variables $\atge{w}{v}$, $\atge{w}{v'}$ on the same
  view variable $W$ with $v<v'$, if $\Def(\atge{w}{v})$ and
  $\Def(\atge{w}{v'})$ appear in $F_t$, then either a bound variable is defined
  for some value strictly between $v$ and $v'$, or the constraint
  $\atge{w}{v'}\Rightarrow\atge{w}{v}\geq 1$ appears in $F_t$.

\item[\inv{Bound}: boundary pins.]
  If any $\atge{w}{v}$ with $v\leq l_W$ is defined, then for the largest such
  $v$ the unit $\atge{w}{v}$ appears in $F_t$; if any $\atge{w}{v}$ with
  $v>u_W$ is defined, then for the smallest such $v$ the unit $\natge{w}{v}$
  appears in $F_t$.  One pin a side is enough: every other out-of-range cut
  follows from it by \invChain{}, through \cref{lem:chain}, which is why the
  arguments that use \inv{Bound} assume \invChain{} as well.  A view variable
  whose bounds are not a
  consequence of the model --- a free bit-sum introduced inside the proof ---
  is exempt, and its owner derives the bounds explicitly instead; no interval
  literal is defined over such a variable today.

\item[\invCut{}: cuts exist.]
  Whenever $\ateq{w}{v}$ is defined, so are $\atge{w}{v}$ and $\atge{w}{v+1}$;
  whenever $\atin{w}{a}{b}$ is defined, so are $\atge{w}{a}$ and
  $\atge{w}{b+1}$.

\item[\invPart{}: the partition.]
  If any interval literal with non-empty in-bounds part is defined for $W$, then
  there is a finite set of \emph{boundaries} $\Bdry(W)\subseteq[l_W,u_W+1]$ with
  $l_W,u_W+1\in\Bdry(W)$.  The \emph{cells} $\Cells(W)$ are the intervals
  $[p,p'-1]$ for consecutive $p<p'$ in $\Bdry(W)$; they partition $[l_W,u_W]$.
  For every cell $C$ the atom $\atset{w}{C}$ is defined --- an equality atom if
  $C$ has width one, an interval atom otherwise.  For every defined two-cut atom
  $\atset{w}{E}$ of $W$ whose in-bounds part $E\cap[l_W,u_W]$ is non-empty, both
  endpoints of that part are boundaries; that is, $E\cap[l_W,u_W]$ is a union of
  adjacent cells.  $\Bdry(W)$ only ever grows.

\item[\invCover{}: coverings.]
  \emph{(i) Split coverings.}  For every $G\subseteq[l_W,u_W]$ with
  $\atset{w}{G}$ defined, $G$ a union of adjacent cells but not itself a cell,
  $F_t$ contains a clause
  $\natset{w}{G}\vee\atset{w}{C_1}\vee\cdots\vee\atset{w}{C_k}$ with $k\geq 2$,
  where every $\atset{w}{C_i}$ is defined, each $C_i$ is a union of adjacent
  cells \emph{properly contained} in $G$, and $\bigcup_i C_i = G$.
  \emph{(ii) Overhanging coverings.}  For every defined $\atset{w}{E}$ whose
  in-bounds part $S := E\cap[l_W,u_W]$ is non-empty and satisfies $E\neq S$,
  $F_t$ contains a clause
  $\natset{w}{E}\vee\atset{w}{C_1}\vee\cdots\vee\atset{w}{C_k}$ where every
  $\atset{w}{C_i}$ is defined, each $C_i$ is a union of adjacent cells, and
  $\bigcup_i C_i = S$.
  \emph{(iii) The root covering.}  If $\Bdry(W)$ exists, $F_t$ contains a clause
  $\atset{w}{C_1}\vee\cdots\vee\atset{w}{C_m}\geq 1$ where every
  $\atset{w}{C_i}$ is defined, each $C_i$ is a union of adjacent cells, and
  $\bigcup_i C_i = [l_W,u_W]$.  A view variable
  with no interval literal, or with only out-of-bounds ones, has no partition and
  no root covering, and this clause holds vacuously --- which is what lets
  \cref{lem:covering} and \cref{thm:complete} still say something about a plain
  order/equality variable, the entire setting of the thesis.

\item[\invCont{}: containment.]
  For any two defined two-cut atoms $\atset{w}{E}$ and $\atset{w}{E'}$ of the
  same view variable $W$ with $E\subsetneq E'$, there is a chain
  $E=E_0\subsetneq E_1\subsetneq\cdots\subsetneq E_m=E'$ with every
  $\atset{w}{E_j}$ defined, such that each clause
  $\natset{w}{E_j}\vee\atset{w}{E_{j+1}}$ appears in $F_t$.

\item[\invView{}: view closure and linking.]
  For every $W,W'\in\ViewVars(X)$: if an atomic variable of $W$ is defined, then so is
  the corresponding atomic variable of $W'$ (\cref{def:image} below); and for
  every defined atomic variable of a proper $W$, the two constraints linking it to
  the corresponding atomic variable of the base $X$ appear in $F_t$.

\item[\invName{}: cells are named.]
  Every cell has appeared in some emitted proof line other than its own
  reification.

  Deliberately about cells, not about every interval literal.  A literal whose
  in-bounds part is empty is defined with no partition, no covering and no
  containment edge, and nothing names it; \invView{} does not suffer, because such
  a literal can only arise from an explicit request, and the request path
  establishes \invView{} for itself.  \cref{rem:naming} says why the trigger is
  naming rather than requesting, which is what this invariant is for.

\end{description}

\invChain{} and \inv{Bound} come from the thesis, generalised only by being
stated per view variable.  \invCut{} is implicit there; we make it explicit
because both new literal kinds depend on it.  \invPart{}, \invCover{} and
\invCont{} are the interval apparatus.  \invView{} is the multi-view
apparatus.  \invName{} is a hygiene condition that \invView{} needs for the
interval literals created as cells, and only for those --- see
\cref{rem:naming}.

\begin{remark}
Nothing above depends on the search state.  Every clause named is a tautology of
the encoding, emitted at the top proof level, and derivable there from what
precedes it --- by RUP, except for the ge-links, which are cutting-planes steps
(see the discussion after \cref{eq:inlink}).  The alternative --- emitting a covering over whatever facts currently
witness an exclusion --- is also sound, but makes the invariant depend on what
was live when, and drags a level-aware registry of live conclusions into the
encoding layer.  It was implemented, and rejected for complexity rather than for
soundness.
\end{remark}

\paragraph{Correspondence between view variables.}

\begin{definitionT}[Image of a literal]
\label{def:image}
Let $W,W'\in\ViewVars(X)$ and let $\theta = \varphi_{W'}\circ\psi_W$, an affine
bijection $\mathbb Z\to\mathbb Z$ with $\theta(v) = s v + d$ where
$s = s_W s_{W'}$ and $d = c_{W'} - s\,c_W$.  The \emph{image} on $W'$ of a
literal $\ell$ over an atomic variable of $W$, written $\ell^{W'}$, is the
literal satisfied by exactly the same values of $X$:
\[
  \begin{array}{llll}
    (\atge{w}{v})^{W'} &=
      \begin{cases}
        \atge{w'}{\theta(v)} & s=+1\\[7pt]
        \natge{w'}{\theta(v)+1} & s=-1
      \end{cases}
    &\qquad
    (\ateq{w}{v})^{W'} &= \ateq{w'}{\theta(v)}\\[2.6ex]
    (\atin{w}{a}{b})^{W'} &=
      \begin{cases}
        \atin{w'}{\theta(a)}{\theta(b)} & s=+1\\[2pt]
        \atin{w'}{\theta(b)}{\theta(a)} & s=-1
      \end{cases}
    &\qquad
    (\natin{w}{a}{b})^{W'} &= \overline{(\atin{w}{a}{b})^{W'}}
  \end{array}
\]
Negation commutes with taking images --- the last entry above, and likewise for
the other two kinds --- so the negated image of an atom may be written by barring
the atom inside, which is what \cref{eq:gelink,eq:eqlink,eq:inlink} do.  Taking
images is functorial: $(\ell^{W'})^{W''} = \ell^{W''}$, and $\ell^{W}=\ell$.
\end{definitionT}

Two facts about \cref{def:image} are worth stating because they are what make
the design tractable.  First, $\theta$ preserves \emph{width}: an interval of
width $k$ maps to an interval of width $k$, so the width-one identification of
\cref{def:widthone} is stable, and the interval machinery never has to consider
a width-one case on either side.  Second, a sign flip converts a positive bound
literal into a negative one, but leaves the polarity of equality and interval
literals alone --- an interval is order-agnostic, a half-line is not.

The atom-level linking constraints are exactly the clausal form of the biconditional
$\ell \Leftrightarrow \ell^{X}$, for each defined atom of each proper
view variable:
\begin{align}
  \text{ge-link:}\quad
    &\natge{w}{v}\vee (\atge{w}{v})^{X}, \qquad
     (\natge{w}{v})^{X}\vee \atge{w}{v};
     \label{eq:gelink}\\
  \text{eq-link:}\quad
    &\nateq{w}{v}\vee \ateq{x}{\psi_W(v)}, \qquad
     \nateq{x}{\psi_W(v)}\vee \ateq{w}{v};
     \label{eq:eqlink}\\
  \text{in-link:}\quad
    &\natin{w}{a}{b}\vee (\atin{w}{a}{b})^{X}, \qquad
     (\natin{w}{a}{b})^{X}\vee \atin{w}{a}{b}.
     \label{eq:inlink}
\end{align}

The eq-link and the in-link are RUP at the point they are emitted.  \textbf{The
ge-link is not}, and this is worth being precise about, because the natural
assumption is that it is.  Summing $\Def(\atge{w}{v})$, $\ViewLink(W)$ and the
corresponding half of $\Def((\atge{w}{v})^X)$ does make the $\BinEnc$ terms
cancel exactly, and saturating leaves the clause --- but that is a
\emph{cutting-planes} derivation, not a unit-propagation one.  What would make
the configuration RUP is Theorem~2.9 of the thesis, and it requires the middle
constraint's right-hand side to lie in $\{0,1\}$; here that constraint is a half
of $\ViewLink(W)$, so the right-hand side is $\pm c_W$, and the theorem does not
apply once $|c_W|>1$.  Example~2.15 of the thesis is a counterexample of exactly
this shape.

Nor is this a corner case.  Checking the clause as a RUP line with a verifier
over a three-bit base, only $c_W = 0$ succeeds at every threshold; $X\in[0,7]$
with $W = X+1$ and $v=6$ is a concrete rejection, and other offsets fail at some
thresholds and not others.  The solver is unaffected, because
\code{need\_gevar} emits these by \code{pol}; but a future simplification of that
emission to a RUP line would fail on the first view with a nonzero offset, and
would fail intermittently, which is worse.

For the eq-link, asserting one side and the negation of the other propagates both
cuts through \cref{eq:deq1}, carries each across by the ge-links, and closes on
\cref{eq:deq2} of the far side.  For the in-link, the same argument with
\cref{eq:din1,eq:din2}, noting that for $s=-1$ the two cuts exchange roles ---
which is exactly the endpoint swap in \cref{def:image}.

\begin{remarkN}[Both in-link clauses are needed, and neither for its implication]
\label{rem:bothlinks}
The two clauses of \cref{eq:inlink} are, as \emph{implications},
$\atin{w}{a}{b}\rightarrow(\atin{w}{a}{b})^X$ and its converse, and both are
already derivable by unit propagation from
\cref{eq:din1,eq:din2} and the ge-links --- one forward
reification, two ge-links, one reverse reification.  That observation is
correct and it is beside the point.  Unit propagation uses a clause in whichever
direction has a unit, and the directions that matter here are the
contrapositives,
\[
  (\natin{w}{a}{b})^{X} \implies \natin{w}{a}{b}
  \qquad\text{and}\qquad
  \natin{w}{a}{b} \implies (\natin{w}{a}{b})^{X},
\]
and neither is derivable without its clause.  A negated interval literal is a
unit on \emph{neither} of its cuts, by the asymmetry noted after
\cref{eq:din2}; all it can do is descend its own view variable's containment
edges, and it bottoms out at cells that are themselves unlinked interval
literals unless they happen to have width one.  Descending to width one is
precisely the per-value shattering that interval literals exist to avoid.  The
underlying error is to decompose a negated interval as if it were a conjunction
of two bounds; it is a \emph{disjunction} of two bounds, and unit propagation
cannot carry a disjunction across one disjunct at a time.

By contrast, one clause of each of \cref{eq:gelink,eq:eqlink} would be enough
for the positive direction, and the pair is emitted for the negative one in
exactly the same way; the interval case is not special in needing two clauses,
only in that the shortcut is tempting.
\end{remarkN}

\paragraph{Defined domains over a view family.}
Any set of atomic literals $L$ defined over view variables of the same base
$X$ represents a restriction of $X$ to some feasible set of values.  We call
the values allowed the \emph{defined domain}, stated on the base's value scale,
so that literals over different view variables can be intersected directly:
\begin{multline}
  \dm_L(X) := \{k : \forall\,\atge{w}{v}\in L .\ \varphi_W(k)\geq v\}
      \cap \{k : \forall\,\natge{w}{v}\in L .\ \varphi_W(k)< v\} \\
    {}\cap \{k : \forall\,\ateq{w}{v}\in L .\ \varphi_W(k) = v\}
      \cap \{k : \forall\,\nateq{w}{v}\in L .\ \varphi_W(k)\neq v\} \\
    {}\cap \{k : \forall\,\atin{w}{a}{b}\in L .\ a\leq\varphi_W(k)\leq b\}
      \cap \{k : \forall\,\natin{w}{a}{b}\in L .\
        \varphi_W(k)<a \text{ or } \varphi_W(k)>b\}.
  \label{eq:domL}
\end{multline}
This set may be empty, or infinite, since $\dm_{\emptyset}(X)=\mathbb Z$.  When
$\ViewVars(X)=\{X\}$ and no interval literal is defined, \cref{eq:domL} is exactly
the thesis's definition.  We write $\dm_L(W) := \varphi_W(\dm_L(X))$ for the
same set on $W$'s scale.

Conversely, the complete (potentially infinite) set of atomic literals over
\emph{every} view variable of $X$ that should be true for a domain
$\dm_t(X)$ is
\begin{align}
  \ImpLits(\dm_t(X)) := \bigcup_{W\in\ViewVars(X)}\Bigl(
    &\{\ateq{w}{v} : \dm_t(W)=\{v\}\} \notag\\
    {}\cup{}&\{\nateq{w}{v} : v\notin\dm_t(W)\} \notag\\
    {}\cup{}&\{\atge{w}{v} : \forall\,v'<v .\ v'\notin\dm_t(W)\} \notag\\
    {}\cup{}&\{\natge{w}{v} : \forall\,v'\geq v .\ v'\notin\dm_t(W)\} \notag\\
    {}\cup{}&\{\atin{w}{a}{b} : \dm_t(W)\subseteq[a,b]\} \notag\\
    {}\cup{}&\{\natin{w}{a}{b} : \dm_t(W)\cap[a,b]=\emptyset\}\Bigr).
  \label{eq:implits}
\end{align}

\paragraph{Propagation results, one view variable at a time.}
We first establish everything for a single view variable $W$, writing $L^W$ for
the atomic literals over atomic variables of $W$ defined in $F_t$, and $[l,u]$
for $[l_W,u_W]$.  \cref{lem:transport} then lifts the results to a family.

It is convenient to normalise away positive two-cut literals once and for all.

\begin{lemmaT}[Chain propagation]
\label{lem:chain}
If \invChain{} is respected, then for any $W$, if $\atge{w}{v}$ and
$\atge{w}{v'}$ are both defined in $F_t$ for $v>v'$, then if $\atge{w}{v}$
propagates so does $\atge{w}{v'}$, and if $\natge{w}{v'}$ propagates
so does $\natge{w}{v}$.
\end{lemmaT}

\begin{proof}
Let $v = v_1 > v_2 > \cdots > v_k = v'$ be the finite sequence of values for
which $\atge{w}{v_j}$ is defined, with nothing defined strictly between
consecutive terms.  Since only finitely many bound literals are defined in a
finite formula, this sequence exists.  \invChain{} puts
$\atge{w}{v_j}\Rightarrow\atge{w}{v_{j+1}}\geq 1$ in $F_t$ for each $j$, and each
is a binary clause, so it propagates in one direction from $\atge{w}{v_j}$ and
in the other from $\natge{w}{v_{j+1}}$.  Chaining gives both claims.
\end{proof}

\begin{lemmaN}[Cut normalisation]
\label{lem:normalise}
Let $R\subseteq L^W$ and let $R'$ be obtained from $R$ by replacing each
positive literal $\ateq{w}{v}$ by $\{\atge{w}{v},\natge{w}{v+1}\}$ and
each positive literal $\atin{w}{a}{b}$ by
$\{\atge{w}{a},\natge{w}{b+1}\}$.  Then every literal of $R'$
propagates under $R$, and $\dm_{R'}(W)=\dm_R(W)$.
\end{lemmaN}

\begin{proof}
The replacements are exactly what \cref{eq:deq1,eq:din1} propagate from the
positive literal, and the cuts are defined by \invCut{}.  The two literal sets
allow the same values by \cref{eq:defeq,eq:defin}.
\end{proof}

Call a set of literals \emph{cut-normalised} if it contains no positive equality
or interval literal.  \cref{lem:normalise} says that any set may be replaced by
a cut-normalised one with the same defined domain, all of whose literals still
propagate, so in the proofs below we may assume that the only two-cut literals
present are negative.  This is where the asymmetry bites: normalisation cannot be applied to
a negative two-cut literal, and there is nothing else to do with one except
propagate through the machinery of \invCover{} and \invCont{}.

\begin{lemmaN}[Containment closure]
\label{lem:contclosure}
Suppose that whenever a two-cut atom $\atset{w}{E}$ of $W$ is defined, the clause
$\natset{w}{E}\vee\atset{w}{E'}$ is emitted for every \emph{minimal}
$E'\supsetneq E$ whose atom is defined, and the clause
$\natset{w}{E''}\vee\atset{w}{E}$ for every \emph{maximal} $E''\subsetneq E$
whose atom is defined.  Then \invCont{} holds.
\end{lemmaN}

\begin{proof}
By induction on the number of defined atoms.  Let $\atset{w}{E}$ and
$\atset{w}{E'}$ both be defined with $E\subsetneq E'$, and consider the later of
the two to be defined, say at step $n$; by the inductive hypothesis \invCont{}
held over the atoms defined before step $n$.

If $\atset{w}{E'}$ is the later, then $\atset{w}{E}$ was already defined, so $E$
is contained in some maximal $E''\subsetneq E'$ with $E\subseteq E''$ and
$\atset{w}{E''}$ defined.  If $E=E''$ the edge
$\natset{w}{E}\vee\atset{w}{E'}$ is emitted directly; otherwise $E\subsetneq
E''$, so the inductive hypothesis gives a chain from $E$ to $E''$, which the
emitted edge $\natset{w}{E''}\vee\atset{w}{E'}$ extends.  The case where
$\atset{w}{E}$ is the later is symmetric, using a minimal container.  Emitting
an edge never invalidates a chain, so \invCont{} continues to hold.
\end{proof}

\begin{remark}
Interval requests may overlap without nesting --- $[7,15]$ after $[5,10]$ ---
so the containment relation is a DAG, not a tree, and a set may have several
incomparable minimal containers.  \cref{lem:contclosure} does not care.  Note
also that it needs \emph{minimal} and \emph{maximal} to be computed over all
sets whose atoms are currently defined for $W$, not merely over the cells.
\end{remark}

\begin{lemmaN}[Cell exclusion]
\label{lem:cell}
Assume \invChain{}, \invCut{}, \invPart{} and \invCont{}.  Let $R$ be a set of
literals, let $U\subseteq L^W$ be a cut-normalised set of literals each of which
propagates under $R$, suppose $F_t\!\restriction_R$ does not propagate to conflict, and
write $D=\dm_U(W)$.  Let $C\in\Cells(W)$ be a cell with $D\cap C=\emptyset$.
Then $\natset{w}{C}$ propagates under $R$.
\end{lemmaN}

\begin{proof}
Write $C=[p,q]$, so that $\atset{w}{C}$ is $\atin{w}{p}{q}$ if $p<q$ and
$\ateq{w}{p}$ if $p=q$; in either case \cref{eq:deq1,eq:din1} give the single PB
constraint $\atset{w}{C}\Rightarrow\atge{w}{p}+\natge{w}{q+1}\geq 2$, and both
cuts are defined by \invCut{}.  That constraint is not a pair of clauses, but it
propagates like one in each direction: $\atset{w}{C}$ makes both cuts units, and
either $\natge{w}{p}$ or $\atge{w}{q+1}$ propagates $\natset{w}{C}$, which is
the direction used below.

\emph{Case 1: some negative two-cut literal of $U$ names a set containing $C$.}
That is, $\natset{w}{E}\in U$ for a defined $\atset{w}{E}$ with $C\subseteq E$.
If $E=C$ we are done.  Otherwise $C\subsetneq E$, and \invCont{} gives a chain of
binary clauses $\natset{w}{C}\vee\atset{w}{E_1}$,
$\natset{w}{E_1}\vee\atset{w}{E_2}$, \dots,
$\natset{w}{E_{m-1}}\vee\atset{w}{E}$; reading each contrapositive in turn from
$\natset{w}{E}$ propagates $\natset{w}{C}$.

\emph{Case 2: no negative two-cut literal of $U$ names a set containing $C$.}
We claim that then no negative two-cut literal of $U$ names a set meeting $C$ at
all.  Let $\natset{w}{E}\in U$.  If $E\cap[l,u]=\emptyset$ then
$E\cap C=\emptyset$ since $C\subseteq[l,u]$.
Otherwise, by \invPart{}, $E\cap[l,u]$ is a union of adjacent cells, and $C$ is
a cell, so either $C\subseteq E$ --- excluded by the case assumption --- or
$C\cap E=\emptyset$.

So every value of $C$ is excluded by a bound literal of $U$.  Put
$v^{+}=\max\{v : \atge{w}{v}\in U\}$ (with $v^{+}=-\infty$ if there is none)
and $v^{-}=\min\{v : \natge{w}{v}\in U\}$ (with $v^{-}=+\infty$ if
there is none).  The values \emph{not} excluded by bounds of $U$ are exactly
$[v^{+},v^{-}-1]$.  If $v^{-}\leq v^{+}$ then $\atge{w}{v^{+}}$ and
$\natge{w}{v^{-}}$ propagate and \cref{lem:contra-bounds} gives a
conflict, contrary to assumption.  So $[v^{+},v^{-}-1]$ is a non-empty interval
disjoint from the interval $C$, and $C$ lies wholly on one side of it.
\begin{itemize}[itemsep=1pt]
  \item If $q<v^{+}$: then $q+1\leq v^{+}$, and $\atge{w}{q+1}$ is defined, so
    \cref{lem:chain} propagates it from $\atge{w}{v^{+}}$, and the constraint
    above propagates $\natset{w}{C}$.
  \item If $p\geq v^{-}$: then $\atge{w}{p}$ is defined and \cref{lem:chain}
    propagates $\natge{w}{p}$ from $\natge{w}{v^{-}}$, and the constraint
    above propagates $\natset{w}{C}$.
    \qedhere
\end{itemize}
\end{proof}

\begin{lemmaN}[Covering completeness]
\label{lem:covering}
Assume \invChain{}, \inv{Bound}, \invCut{}, \invPart{} and \invCover{}.  Let
$\atset{w}{E}$ be a defined two-cut atom of $W$ and let $S=E\cap[l,u]$.
\begin{enumerate}[label=(\alph*),itemsep=1pt]
  \item If $S=\emptyset$ then $\natset{w}{E}$ propagates under any $R$.
  \item If $S\neq\emptyset$, $E$ is not currently a cell, and $\natset{w}{C}$
    propagates for every cell $C\subseteq S$, then $\natset{w}{E}$ propagates.
\end{enumerate}
\end{lemmaN}

\begin{proof}
\emph{(a)}  Write $E=[a,b]$.  If $b<l$ then $\atge{w}{b+1}$ is defined by
\invCut{} with $b+1\leq l$, so \inv{Bound} makes $\atge{w}{v}$ a unit for some
$b+1\leq v\leq l$, \cref{lem:chain} propagates $\atge{w}{b+1}$ from it, and
\cref{eq:deq1,eq:din1} propagate $\natset{w}{E}$.  If $a>u$ then $\atge{w}{a}$
is defined with $a>u$, so \inv{Bound} makes $\natge{w}{v}$ a unit for some
$u<v\leq a$, \cref{lem:chain} propagates $\natge{w}{a}$ from it, and
\cref{eq:deq1,eq:din1} propagate $\natset{w}{E}$ through the other cut.

\emph{(b)}  For a set that is a union of adjacent cells, define its \emph{rank}
to be the number of cells it contains.  We first show, by induction on rank,
that $\natset{w}{G}$ propagates for every $G\subseteq S$ with $\atset{w}{G}$
defined and $G$ a union of adjacent cells.

If $G$ has rank one then $G$ is a cell, since its two endpoints are boundaries
and no boundary lies strictly inside it; as $G\subseteq S$, $\natset{w}{G}$
propagates by hypothesis.  If $G$ has rank at least two then $G$ is not a cell,
so \invCover{}(i) puts
$\natset{w}{G}\vee\atset{w}{C_1}\vee\cdots\vee\atset{w}{C_k}$ in $F_t$ with
every $\atset{w}{C_i}$ defined and each $C_i$ a union of adjacent cells,
properly contained in $G$ --- hence of strictly smaller rank --- and contained
in $S$.  By induction each $\natset{w}{C_i}$ propagates, so the covering
propagates $\natset{w}{G}$.

Now for $E$ itself.  If $E\subseteq[l,u]$ then $E=S$, which is a union of
adjacent cells by \invPart{}, and the claim applies directly.  Otherwise
$E\neq S$, and \invCover{}(ii) puts
$\natset{w}{E}\vee\atset{w}{C_1}\vee\cdots\vee\atset{w}{C_k}$ in $F_t$ with
every $\atset{w}{C_i}$ defined, each $C_i$ a union of adjacent cells, and
$\bigcup_i C_i=S$; each $C_i\subseteq S$, so $\natset{w}{C_i}$ propagates by the
claim, and the covering propagates $\natset{w}{E}$.
\end{proof}

\begin{remark}
\invCover{} is stated structurally, over the \emph{current} partition, but an
implementation emits each clause once and never revisits it: a split covering
when a cell is first split into two halves, and a request covering when a
literal spanning several cells is defined.  Both statements are then maintained
by later refinement without re-emission, which is the whole content of
``coverings compose across refinements'' --- a covering whose pieces are later
subdivided still satisfies \invCover{}, because the pieces remain defined
literals and remain unions of adjacent cells.  This is why \cref{lem:covering}
is proved by induction on rank rather than by a direct appeal to the clauses
that happen to have been written: the induction descends through whatever
coverings exist now, not through the order in which they were emitted.

The distinction matters for one specific reason.  An equality atom created
before its view variable has any partition is not covered by any clause at
that point; when the partition is later created it becomes a width-one cell,
and a width-one cell can never be split, so it never needs a covering.  A
formulation of \invCover{} in terms of ``was a cell when defined'' would have to
special-case it; the structural form does not.
\end{remark}

\begin{theoremT}[Empty defined domain propagates to conflict]
\label{thm:wipeout}
Let $F_t$ be a set of PB constraints either in the encoding of a CP problem, as
produced by \cref{proc:encoding}, or derived from it, and suppose \invChain{}
and \invCut{} are respected.  Let $W$ be a view variable, and let $U\subseteq
L^W$ be a set of literals each of which propagates under some $R$.  If
$\dm_U(W)=\emptyset$ then $F_t\!\restriction_R$ propagates to conflict.
\end{theoremT}

\begin{proof}
Suppose not, and cut-normalise $U$ using \cref{lem:normalise}, which preserves
both the defined domain and the property of propagating under $R$.

Since $\dm_U(W)=\emptyset$, every integer is excluded by $U$ somehow.  In
particular the values below any fixed point must be excluded.  No literal
$\natge{w}{v}$ excludes anything below $v$, and $U$ is finite --- it is a subset
of $L^W$, and only finitely many atomic variables are defined in a finite
formula --- so its negative two-cut literals between them exclude only a bounded
set.  Hence $U$ contains at least one literal $\atge{w}{v}$.  Let $v^{+}$ be the
largest value such that $\atge{w}{v^{+}}$ propagates under $R$; this exists for
the same reason.
Symmetrically, let $v^{-}$ be the smallest value such that
$\natge{w}{v^{-}}$ propagates; this exists for the mirrored reason.

If $v^{+}\geq v^{-}$ then \cref{lem:contra-bounds} gives a conflict.  So
$v^{+}<v^{-}$, and in particular the value $v^{+}$ is not excluded by any bound
literal of $U$.  It must therefore be excluded by a negative two-cut literal of
$U$, and there are two cases.
\begin{itemize}[itemsep=1pt]
  \item $\nateq{w}{v^{+}}\in U$.  Then $\atge{w}{v^{+}}$ propagates,
    and \cref{eq:deq2} --- the clause
    $\ateq{w}{v^{+}}\vee\natge{w}{v^{+}}\vee\atge{w}{v^{+}+1}$ ---
    propagates $\atge{w}{v^{+}+1}$, contradicting the maximality of $v^{+}$.
  \item $\natin{w}{a}{b}\in U$ with $a\leq v^{+}\leq b$.  By
    \invCut{}, $\atge{w}{a}$ is defined, and $a\leq v^{+}$, so
    \cref{lem:chain} propagates it.  Then \cref{eq:din2} --- the clause
    $\atin{w}{a}{b}\vee\natge{w}{a}\vee\atge{w}{b+1}$ --- propagates
    $\atge{w}{b+1}$, and $b+1>v^{+}$, again contradicting maximality.
    \qedhere
\end{itemize}
\end{proof}

\begin{remarkN}[The root covering is not load-bearing here]
\label{rem:rootcovering}
\cref{thm:wipeout} uses only the reverse reifications, the order chain, and the
fact that both cuts of a two-cut literal are defined.  It does \emph{not} use
\invCover{} at all, and in particular not the root covering \invCover{}(iii),
which is the clause one would naively expect to be responsible for detecting
wipeout at the literal level.  The reason is the second bullet above: a negated
interval literal, combined with a bound that has reached its lower cut, jumps
the bound over the whole interval in one step.  Walking those jumps from the
bottom of the domain is exactly a traversal of the excluded region.

The root covering is nonetheless emitted, and we do not propose removing it.  It
is one clause per variable; it is the flag-level analogue of the direct
encoding's at-least-one clause; and the argument above is a hand analysis of
unit propagation, which is precisely the kind of argument that has been wrong
before in this design --- see \nsec{sec:loadbearing}.
\end{remarkN}

\begin{theoremT}[Complete propagation of implied atomic literals]
\label{thm:complete}
Let $F_t$ be as in \cref{thm:wipeout}, and suppose \invChain{}, \inv{Bound},
\invCut{}, \invPart{}, \invCover{} and \invCont{} are respected.  Let $W$ be a
view variable and $U\subseteq L^W$ a set of literals each of which propagates
under some $R$.  If $F_t\!\restriction_R$ does not propagate to contradiction,
then every literal in $\ImpLits(\dm_U(X))\cap L^W$ propagates under $R$.
\end{theoremT}

\begin{proof}
Cut-normalise $U$ by \cref{lem:normalise}, and write $D=\dm_U(W)$, which is
non-empty since otherwise \cref{thm:wipeout} would give a contradiction.
Consider a literal $\ell\in\ImpLits(\dm_U(X))\cap L^W$.  By \cref{eq:implits}
there are six cases.

\emph{Case 1: $\ell=\atge{w}{v}$, with $v'\notin D$ for all $v'<v$.}  In
particular $D$ is bounded from below.  Since $U$ is finite, its negative two-cut
literals exclude only a bounded set, so $U$ must contain a positive bound
literal, and we may let $v^{+}$ be the largest value with $\atge{w}{v^{+}}$
propagating.
If $v^{+}>v$ then $\ell$ propagates by \cref{lem:chain}.  If $v^{+}=v$ then
$\ell$ propagates by definition of $v^{+}$.  Otherwise $v^{+}<v$, so
$v^{+}\notin D$, and $v^{+}$ is not excluded by any bound literal of $U$ (by
maximality of $v^{+}$ and, on the other side, because
$\natge{w}{v'}\in U$ with $v'\leq v^{+}$ would give a conflict by
\cref{lem:contra-bounds}).  So $v^{+}$ is excluded by a negative two-cut
literal, and exactly as in the two bullets of \cref{thm:wipeout} we derive
$\atge{w}{v''}$ for some $v''>v^{+}$, contradicting maximality.

\emph{Case 2: $\ell=\natge{w}{v}$.}  Symmetric to Case 1, taking
$v^{-}$ the smallest value with $\natge{w}{v^{-}}$ propagating and
using the other cut in each reverse reification.

\emph{Case 3: $\ell=\ateq{w}{v}$, so $D=\{v\}$.}  Since $\ell$ is defined, so
are $\atge{w}{v}$ and $\atge{w}{v+1}$ by \invCut{}; both
$\atge{w}{v}\in\ImpLits$ and $\natge{w}{v+1}\in\ImpLits$, so they
propagate by Cases 1 and 2, and \cref{eq:deq2} then propagates $\ell$.

\emph{Case 4: $\ell=\nateq{w}{v}$, so $v\notin D$.}  Consider what excludes $v$;
after cut normalisation the candidates are a bound literal and a negative
two-cut literal.  If a bound literal does, then either $\atge{w}{v+1}$ or
$\natge{w}{v}$ propagates by \cref{lem:chain}, and \cref{eq:deq1} then
propagates $\ell$ from whichever it is.  If instead a negative two-cut
literal $\natset{w}{E}\in U$ with $v\in E$ does, then $\{v\}\subseteq E$; if
$E=\{v\}$ then $\ell\in U$, and otherwise \invCont{} gives a chain of binary
clauses whose contrapositives propagate $\ell$ from $\natset{w}{E}$.  If a
positive two-cut literal excludes $v$, cut normalisation has already replaced it
by bound literals, so this is the first case.

\emph{Case 5: $\ell=\atin{w}{a}{b}$, with $D\subseteq[a,b]$.}  Since $\ell$ is
defined, $\atge{w}{a}$ and $\atge{w}{b+1}$ are defined by \invCut{}.  Every
$v<a$ is outside $D$ and every $v\geq b+1$ is outside $D$, so
$\atge{w}{a}$ and $\natge{w}{b+1}$ are in
$\ImpLits(\dm_U(X))\cap L^W$ and propagate by Cases 1 and 2.  \cref{eq:din2}
then propagates $\ell$.

\emph{Case 6: $\ell=\natin{w}{a}{b}$, with $D\cap[a,b]=\emptyset$.}
Write $E=[a,b]$ and $S=E\cap[l,u]$.  If $S=\emptyset$ then $\ell$ propagates by
\cref{lem:covering}(a).  Otherwise, $S$ is a union of cells by \invPart{}, and
every cell $C\subseteq S$ satisfies $D\cap C=\emptyset$, so $\natset{w}{C}$
propagates by \cref{lem:cell}.  If $E$ is itself a current cell we are done; if
not, \cref{lem:covering}(b) propagates $\ell$.
\end{proof}

\begin{remark}[What ``complete'' does and does not mean]
\label{rem:domlimits}
\cref{eq:domL} takes no account of $\Bnd(W)$, so $\dm_U(W)$ can contain values
the variable's own bounds exclude, and \cref{thm:complete} is correspondingly
narrower than its name suggests.  With cells $[1,3]$, $[4,6]$, $[7,10]$ on
$[1,10]$ and $U = \{\natset{w}{[1,3]}, \natset{w}{[4,6]}\}$, unit propagation does
derive $\atge{w}{7}$ --- through the root covering and the surviving cell --- but
$\atge{w}{7}\notin\ImpLits(\dm_U)$, so the theorem does not claim it.  With
$U$ extended by $\natset{w}{[7,10]}$, unit propagation reaches a conflict while
$\dm_U(W)$ is still non-empty, so \cref{thm:wipeout}'s hypothesis does not hold
either.  Neither is a counterexample to anything stated; both are consequences of
$\dm_U$ being computed from $U$ alone.

This is the thesis's convention and we keep it, because it is what makes the
theorems compose: \cref{sec:unsat} puts the root bound literals into the
propagated set, at which point $\dm_U(W)\subseteq[l_W,u_W]$ and the gap closes.
The reader should simply not read ``complete propagation'' as ``everything the
solver's domain implies''.
\end{remark}

\paragraph{Lifting to a view family.}

\begin{lemmaN}[Unit transport]
\label{lem:transport}
Assume \invView{}.  Let $W,W'\in\ViewVars(X)$ and let $\ell$ be a literal over an
atomic variable of $W$ defined in $F_t$.  If $\ell$ propagates under $R$, then
$\ell^{W'}$ propagates under $R$.
\end{lemmaN}

\begin{proof}
It suffices to show that a literal crosses a single link pair, in either
direction and either polarity; a general $W\to W'$ transport is then at most two
such crossings, through the base, because images compose:
$(\ell^{X})^{W'} = \ell^{W'}$.  So take $W$ proper and consider its link to the
base.  By \invView{} the atomic variable underlying $\ell^{X}$ is defined and the
two linking clauses
\cref{eq:gelink}, \cref{eq:eqlink} or \cref{eq:inlink} appear in $F_t$.  Those
two clauses are $\overline{p}\vee q$ and $\overline{q}\vee p$ for
$p$ the positive literal of $W$'s atom and $q$ its image; a binary clause
propagates its other literal from a unit, so $p\mapsto q$ uses the first and
$\overline{q}\mapsto\overline{p}$ the same one, while $q\mapsto p$ and
$\overline{p}\mapsto\overline{q}$ use the second.  All four polarity
combinations are therefore covered.
\end{proof}

\begin{remarkN}[Why naming, not requesting, is the trigger]
\label{rem:naming}
\cref{lem:transport} is only as good as \invView{}, and \invView{} is the
invariant that is easiest to break silently.  Inequality and equality atoms are
mirrored eagerly at the moment they are defined, so for them \invView{} is
maintained by construction.  Interval literals are not, because they come into
existence two ways: because some consumer asked for one --- a conclusion, a
reason element, a branching decision --- or because the partition machinery
created it as a cell, which no consumer ever asked for.  A rule that mirrors
``requested'' literals leaves every cell unlinked, and a later decision, reason
or conclusion that happens to name one of those cells then has no link at all:
a unit-propagation dead end, with nothing wrong at the point of failure.

The correct trigger is \emph{naming}: an interval literal is mirrored and linked
the first time it appears in an emitted proof line, from whichever side names it
first.  \invName{} then says that this covers every defined interval literal,
and it does, because every creation path writes a line naming the literal it has
just created, in the same call: a split emits its split covering, a partition
creation its root covering, a request its own covering, and a definition its
containment edges.  Only the \code{red} reification lines name a literal without
going through the naming hook.  The invariant to preserve is therefore: every
non-\code{red} emission names its literals.  A future code path that rendered an
interval literal into a line without going through the naming hook would leave
one unlinked, and nothing would notice.
\end{remarkN}

\setcounter{theoremP}{1}
\begin{theoremP}[Empty defined domain, family version]
\label{thm:wipeout-family}
Let $F_t$ be as in \cref{thm:wipeout}, and suppose \invChain{}, \invCut{} and
\invView{} are respected.  Let $L$ be the atomic literals defined in $F_t$ for
any view variable of $X$, and let $U\subseteq L$ be a set of literals each of
which propagates under $R$.  If $\dm_U(X)=\emptyset$ then
$F_t\!\restriction_R$ propagates to conflict.
\end{theoremP}

\begin{proof}
Fix any $W\in\ViewVars(X)$ and let $U^W = \{\ell^{W} : \ell\in U\}$.  By
\cref{lem:transport} every literal of $U^{W}$ propagates under $R$, and
$U^W\subseteq L^W$ by \invView{}.  Images preserve satisfaction by base values,
so $\dm_{U^W}(X)=\dm_U(X)=\emptyset$.  \cref{thm:wipeout} applied to $U^W$ gives
the conflict.
\end{proof}

\begin{theoremP}[Complete propagation, family version]
\label{thm:complete-family}
Let $F_t$, $L$ and $U$ be as in \cref{thm:wipeout-family}, and suppose in
addition that \inv{Bound}, \invPart{}, \invCover{} and \invCont{} are respected
for every view variable of $X$.  If $F_t\!\restriction_R$ does not propagate to
contradiction, then every literal in $\ImpLits(\dm_U(X))\cap L$ propagates
under $R$.
\end{theoremP}

\begin{proof}
Let $\ell\in\ImpLits(\dm_U(X))\cap L$ be a literal over the view variable $W$.
Transport $U$ to $U^{W}$ as in the previous proof; every literal of $U^W$
propagates and $\dm_{U^W}(X)=\dm_U(X)$.  Then \cref{thm:complete} applied on
$W$ gives that $\ell$ propagates.
\end{proof}

The shape of this argument is worth stating explicitly, because it is what makes
the design maintainable: \textbf{transport, then argue on one side}.  Two
view variables of one variable end up holding the same literals, by \invView{},
but there is no reason for them to hold those literals in the same \emph{roles},
and in practice they do not always: we have observed proofs in which an interval
is a partition cell on one view variable and a request with its own covering on
its image, and in which the two containment DAGs are not isomorphic.
\cref{thm:complete-family} does not care, because \cref{thm:complete} is proved
per view variable from invariants that each side satisfies independently.  That
is deliberate, and it is worth being explicit about why: an argument that
appealed to the two sides having corresponding clause sets would be unsound
reasoning \emph{even in the runs where the correspondence happens to hold},
because nothing maintains it.

\subsubsection*{How facts move: a summary}
\phantomsection\addcontentsline{toc}{subsubsection}{How facts move: a summary}\label{sec:crossings}

\cref{thm:complete-family} says that everything implied propagates.  It is
worth setting out separately \emph{which} clause family carries each kind of
crossing, because that is the level at which the design is modified, and
because a clause family with no crossing to its name is a candidate for
deletion --- correctly or, as has happened three times, incorrectly.

\begin{center}
\small
\begin{tabular}{@{}p{3.9cm}p{3.0cm}p{7.0cm}@{}}
\toprule
\textbf{Fact} & \textbf{Becomes} & \textbf{Carried by} \\
\midrule
\multicolumn{3}{@{}l}{\emph{Within one view variable}}\\[2pt]
$\atge{w}{v}$ & $\atge{w}{v'}$, $v'<v$ & order chain, \invChain{}
  (\cref{lem:chain}) \\
$\natge{w}{v}$ & $\natge{w}{v'}$, $v'>v$ & the same clauses,
  read the other way \\
$\ateq{w}{v}$ or $\atin{w}{a}{b}$ & its two cuts & $\Def^{\Rightarrow}$,
  \cref{eq:deq1,eq:din1} \\
both cuts & $\ateq{w}{v}$ or $\atin{w}{a}{b}$ & $\Def^{\Leftarrow}$,
  \cref{eq:deq2,eq:din2} \\
one cut crossing the literal & $\nateq{w}{v}$ or
  $\natin{w}{a}{b}$ & $\Def^{\Rightarrow}$, contrapositive \\
$\nateq{w}{v}$ or $\natin{w}{a}{b}$ & a bound stepping over
  it & $\Def^{\Leftarrow}$: this is the hole jump, and the reason
  \cref{thm:wipeout} needs no covering \\
$\natset{w}{E'}$, $E\subsetneq E'$ & $\natset{w}{E}$ & containment, \invCont{}
  (\cref{lem:contclosure}) \\
$\natset{w}{C_1},\dots,\natset{w}{C_k}$ & $\natset{w}{E}$ & covering,
  \invCover{} (\cref{lem:covering}) \\
$\atset{w}{E}$ and all but one piece excluded & the surviving piece & covering \\
every cell excluded & conflict & reverse reifications and the chain
  (\cref{thm:wipeout}); the root covering is redundant here, see
  \cref{rem:rootcovering} \\[4pt]
\multicolumn{3}{@{}l}{\emph{Between two view variables of the same variable}}\\[2pt]
$\atge{w}{v}$, either polarity & its image on $X$ & ge-links,
  \cref{eq:gelink}; polarity flips when $s_W=-1$ \\
$\ateq{w}{v}$, either polarity & its image on $X$ & eq-links, \cref{eq:eqlink} \\
$\atin{w}{a}{b}$, either polarity & its image on $X$ & in-links,
  \cref{eq:inlink}; \emph{both} clauses needed, see \cref{rem:bothlinks} \\
every bit of $X$ fixed & every bit of $W$ fixed & $\ViewLink(W)$ and
  thesis Theorem~2.8; the crossing the completion argument of
  \cref{sec:optimality} needs.  Partial assignments cross too --- for $W=X+8$
  over a three-bit base, $\ViewLink(W)$ forces $w_3$ from the empty assignment
  --- but an atomic literal is not carried across this way in general; see
  \nsec{sec:notprovided} \\
anything on $W_1$ & its image on $W_2$ & two hops through the base
  (\cref{lem:transport}) \\[4pt]
\multicolumn{3}{@{}l}{\emph{Between different variables}}\\[2pt]
a value & a value, across $\BinEnc$ equality & thesis Theorem~2.8 \\
opposing bounds & conflict, across $\BinEnc$ equality & thesis Theorems~2.7,~2.9 \\
a lone bound & the same bound on the other side & \textbf{not available in
  general}; see \nsec{sec:notprovided} \\
\bottomrule
\end{tabular}
\end{center}

\subsubsection*{What this does not provide}
\phantomsection\addcontentsline{toc}{subsubsection}{What this does not provide}\label{sec:notprovided}

Two limits are worth stating plainly, because misreading either of them has
cost time.

\paragraph{A lone bound does not cross a constraint.}
The results above are entirely about one CP variable and its view variables.
Nothing here lets unit propagation move a bound literal on $X$ to a bound
literal on $Y$ across a constraint such as $\BinEnc(X)=\BinEnc(Y)$: doing so
would require adding two constraints together, which unit propagation never
does.  Example~2.15 of the thesis exhibits an explicitly contradictory triple of
bit-sum constraints that does not unit propagate at all.  This limit is real,
and it is not repaired by anything in this document.

What removes the difficulty is that no replay obligation ever requires it, once
the vocabulary is complete: reasons and conclusions are single-literal facts
over the vocabulary of \S3.3.1, values cross by thesis
Theorem~2.8, and contradictions by Theorems~2.7 and~2.9.  Where a propagator
genuinely needs a bound to cross, it emits explicit bridge lemmas as part of its
justification.  For an interval conclusion pruned across an equality, those are
the pair
\[
  \atge{x}{a}\Rightarrow\atge{y}{a'}
  \qquad\text{and}\qquad
  \atge{y}{b'+1}\Rightarrow\atge{x}{b+1},
\]
each RUP by the Theorem~2.9 configuration --- which does apply here, the middle
constraint being an equality between two bit-sums, so its right-hand side is
zero.  The lemmas' \emph{conclusions} mention only inequality atoms, which is
what lets them be shared by any interval literal with those endpoints; they are
nevertheless emitted under the propagator's reason, which does contain the
interval literal, and are deleted with the scaffolding level rather than
persisting.  This is
\S3.4 material and we do not develop it here; we note only that the interval
vocabulary does not change the situation, because an interval literal asserts
only inequality atoms, never bits.

We record the historical trap: when the literal layer is incomplete, a replay
stalls, and the stalled trail's \emph{next} step is typically a bound crossing.
The failure therefore looks exactly like ``unit propagation cannot thread a
bound across the equality''.  It usually is not.  Twice, that misdiagnosis
produced elaborate and unnecessary machinery --- global cross-variable
coverings, durable top-level lemmas --- for a problem that was a missing
covering clause somewhere else entirely.

\paragraph{Interval literals require a bit-vector representation.}
\cref{eq:defin} reifies against two inequality atoms, which are in turn reified
against $\BinEnc(W)$.  A variable given a purely direct encoding --- one
$0$--$1$ variable per value, with an at-least-one and an at-most-one clause,
and no bit vector --- has no cuts to reify against, and no interval literals are
available for it.  This is a property of the variable, not of views: whether a
view may carry interval literals is decided by whether its \emph{base} has a bit
vector.  Consumers that would state an interval fact fall back to stating it per
value, which is correct but linear in the width.  There is no such fallback
inside the literal layer itself: a request for an interval literal on a
bit-vector-less variable is an error.  One degenerate case of the direct
encoding does have a bit vector and is therefore not affected; see
\cref{sec:open}.

\subsubsection{Unsatisfiability Proofs from Backtracking Search}
\label{sec:unsat}

We can now explain how a CP solver creates a PB proof when it determines, via
decisions, backtracking and constraint propagation, that a problem is
unsatisfiable.  In short: every time the solver backtracks, it immediately logs
a PB constraint encoding, via atomic literals, the negation of the sequence of
decisions made prior to discovering a conflict.  This can be thought of as a
nogood clause from a solving perspective, but we refer to it as a
\emph{backtracking justification} once written in the proof log.

The solver is not restricted to guessing assignments, but we require that each
branching decision can be expressed as a single atomic literal, and that for the
decision variable all assignments are somehow partitioned by the branches.  In
general, if there are $k$ decisions before a particular backtrack, the required
backtracking justification is
\begin{equation}
  B_t := \sum_{i=0}^{k-1}\overline{d}_i \geq 1
  \label{eq:bt}
\end{equation}
where each $d_i$ is an atomic literal representing the branching decision at
level $i$ of the current search path.

The extended vocabulary widens what may appear as a $d_i$ in two ways, and both
are used.  A decision may be an interval literal, so that interval splitting
(``$X\in[a,b]$'' versus ``$X\notin[a,b]$'') is a two-way branch over one atomic
literal rather than a pair of bounds; and a decision may be stated over any
view variable, though in practice a solver branches on the variables it has
state for, which are the base variables.  Neither changes \cref{eq:bt}.

When the solver backtracks from the root decision level, the sequence of
decisions is empty and \cref{eq:bt} becomes $0\geq 1$: the trivial
unsatisfiability conclusion we want to reach.

For these logged constraints to be efficiently checkable, they need to follow
by RUP.  The idea, going back to Elffers et al.\ and stated explicitly by Gocht
et al., is to maintain an invariant so that every backtracking clause does.  For
any base variable $X$, and any point $t$ immediately prior to backtracking or
immediately before any domain change, let $\dm_t(X)$ be the current domain of
$X$ as stored by the solver; let $B_t$ be the backtracking justification for the
current sequence of decisions; and let $F_t$ be the constraints in the input
encoding or derived in the proof up to $t$.  We must have logged enough that:

\begin{description}[font=\normalfont\bfseries,leftmargin=3.2em,style=nextline]
\item[\invConflict{} (thesis Inv2).]
  If the solver has detected contradiction for the current state of domains,
  then $F_t\cup\{\neg B_t\}$ must unit propagate to contradiction.
\item[\invDomain{} (thesis Inv3).]
  For every base variable $X$, if $L$ is the set of atomic literals defined for
  \emph{any view variable of} $X$ that would propagate during unit propagation
  of $F_t\cup\{\neg B_t\}$, then unless contradiction is reached beforehand, we
  must have $\dm_L(X)\subseteq\dm_t(X)$.
\end{description}

\invDomain{} is where the multi-view setting shows up in the search
invariants, and it shows up in exactly the way one would want.  The solver holds
one domain per base variable; a view has no domain of its own, only an affine
window onto its base's.  Correspondingly, $\dm_L(X)$ in \cref{eq:domL} is an
intersection over the literals of \emph{all} view variables, stated on the
base's value scale.  A fact learned about a view therefore constrains the same
object the solver's domain constrains, and the containment
$\dm_L(X)\subseteq\dm_t(X)$ is between two subsets of the same set of integers.

\begin{theoremT}
\label{thm:btrup}
If \invConflict{} and \invDomain{} are respected, the literal-layer invariants
\cref{thm:wipeout-family} needs hold --- \invChain{}, \invCut{} and \invView{} ---
and a complete solver always logs backtracking justifications, then every
backtracking justification $B_t$ follows by RUP.
\end{theoremT}

\begin{proof}
In a standard backtracking and propagation solver there are three cases in which
the solver backtracks.

\emph{Case 1: a propagator has detected infeasibility.}  Then \invConflict{}
implies $B_t$ is RUP.

\emph{Case 2: a propagator has removed all remaining values from the domain of
some base variable $X$.}  Then \invDomain{} implies that unit propagation on
$F_t\cup\{\neg B_t\}$ either obtains a contradiction, or propagates a set of
literals $L$ over view variables of $X$ with $\dm_L(X)=\emptyset$.  In the
latter case \cref{thm:wipeout-family} implies that further unit propagation must
still result in contradiction.

\emph{Case 3: the solver has exhausted all branching possibilities for its
current branching variable $X$.}  Let $\dm_t(X)$ be the domain of $X$ after
propagation at the current decision level.  The branching choices partition
$\dm_t(X)$.  Since this is not a leaf node, backtracking justifications have
already been logged for each of the $k$ child branches, each of the form
$G\Rightarrow \overline{d}_i\geq 1$ where $G$ encodes the decisions above this
level and $d_i$ is the $i$th branching decision on $X$; so each
$\overline{d}_i$ propagates under $F_t\cup\{\neg B_t\}$.

Now suppose $F_t\cup\{\neg B_t\}$ does not propagate to contradiction, and let
$L$ be the set of atomic literals over view variables of $X$ that propagate
under unit propagation to fixpoint.  We must have $\dm_L(X)\neq\emptyset$, since
otherwise \cref{thm:wipeout-family} gives a conflict; so pick $v\in\dm_L(X)$.
By \invDomain{}, $v\in\dm_t(X)$, so some branch decision $d_i$ is satisfied by
$v$, the branches being a partition of $\dm_t(X)$.  But $\overline{d}_i\in L$,
and $v\in\dm_L(X)$ means $v$ satisfies every literal of $L$, so $v$ does not
satisfy $d_i$ --- a contradiction.  Note that $d_i$ may be stated over any
view variable of $X$, and that $\dm_L$ is computed over all of them, so no
transport step is needed here beyond the one inside
\cref{thm:wipeout-family}.
\end{proof}

\paragraph{Maintaining the invariants.}
If the CP solver happened to perform only inferences no stronger than unit
propagation on the input PB formula, \invConflict{} and \invDomain{} would be
maintained trivially.  But the encoding is deliberately not informing the
solving process.  So we require further justification constraints to be logged.

To define these we make use of \emph{reasons}, which for CP proof logging are
just literal sets $R$ used to reify intermediate facts written in the proof log.
We say a reason $R$ is \emph{valid} for a domain state $\dm_t$ at time $t$ if
every $\ell\in R$ is guaranteed to propagate under $F_t\cup\{\neg B_t\}$, if
contradiction is not reached first.

Whenever the solver detects contradiction from the current domain state, it must
derive in the proof a constraint of the form
\begin{equation}
  R \Rightarrow 0\geq 1
  \label{eq:conflictjust}
\end{equation}
where $R$ is a reason valid at the point of conflict; we call this a
\emph{conflict justification}.  Whenever a propagator makes an inference, it
must derive a \emph{propagation justification}
\begin{equation}
  R \Rightarrow \ell \geq 1
  \label{eq:propjust}
\end{equation}
where $\ell$ is an atomic literal describing the inference and $R$ is a reason
valid immediately before the domain change.

Two things about \cref{eq:propjust} are new.  First, $\ell$ may be a negated
interval literal: a propagator that removes a run of values from a domain states
that as one conclusion $\natin{w}{a}{b}$, not as $b-a+1$ conclusions.
Second, $\ell$ and the literals of $R$ need not be stated over the same
view variable, and routinely are not.  A propagator whose operand is a view
concludes over the view; a reason derived from the solver's domains is stated
over whichever view variable the propagator holds; and a branching decision was
taken on the base.  All three meet only through \cref{lem:transport}.

\begin{lemmaT}
\label{lem:reasons}
If $R$ is a finite set of literals defined in $F_t$ with
$R\subseteq\bigcup_i\ImpLits(\dm_t(X_i))$, and \invChain{}, \inv{Bound},
\invCut{}, \invPart{}, \invCover{}, \invCont{}, \invView{} and \invDomain{} are
respected, then $F_t\cup\{\neg B_t\}$ either propagates to contradiction or
propagates every literal in $R$.
\end{lemmaT}

\begin{proof}
Let $\ell\in R$ be a literal over a view variable $W$ of the base variable
$X_i$, so $\ell\in\ImpLits(\dm_t(X_i))$.  Suppose $F_t\cup\{\neg B_t\}$ does not
propagate to contradiction, and let $L$ be the atomic literals over
view variables of $X_i$ that propagate, as guaranteed by \invDomain{}.  Since
we do not reach conflict, $\dm_L(X_i)\neq\emptyset$ by
\cref{thm:wipeout-family}.  Then $\dm_L(X_i)\subseteq\dm_t(X_i)$, so
$\ImpLits(\dm_L(X_i))\supseteq\ImpLits(\dm_t(X_i))$ by \cref{eq:implits}, and
hence $\ell\in\ImpLits(\dm_L(X_i))$.  \cref{thm:complete-family} then tells us
that $\ell$ propagates.
\end{proof}

There are two standard ways to obtain a valid reason.  The atomic literals
representing the current decisions $G$ are always valid for the current subtree,
as they are by definition propagated by $\neg B_t$.  Alternatively, by
\cref{lem:reasons}, any finite set of literals implied by the current domains
will do --- and this is the form that matters here, because it is what lets a
propagator state ``this run of values is absent from the domain'' as a single
reason element rather than as one literal per value.  Our definition of reasons
is more permissive than the reason clauses of lazy clause generation: we retain
the ability to use auxiliary variables from the PB encoding, or fresh variables
introduced by redundance, as part of a reason.

To ensure that justifications suffice to maintain the invariants, we also
require the solver to introduce by redundance the atomic literals for the
initial bounds at the start of the proof, if not already present, and derive by
RUP the units stating that they hold at the root decision level: for a variable
$X$ with initial bounds $l$ and $u$, these are
\begin{equation}
  \atge{x}{l}\geq 1 \qquad\text{and}\qquad \natge{x}{u+1}\geq 1 .
  \label{eq:rootbounds}
\end{equation}
These follow by RUP by the same argument as \cref{cor:chainrup}, since the bound
constraints $\Bnd(X)$ are present, and likewise for every view variable, proper
ones included, because $\Bnd(W)$ is asserted for each.

\begin{remark}[The eager introduction is load-bearing, and the solver is lazier]
\label{rem:lazypins}
\invDomain{} needs \cref{eq:rootbounds} to be present from the start.  With
$L=\emptyset$, \cref{eq:domL} gives $\dm_L(X)=\mathbb Z$, which is not contained
in $\dm_t(X)$, so the base case of \cref{thm:maintain} fails without it.

The implementation does not introduce these eagerly.  Atoms are created only
when a proof line needs them, so a variable no proof line has mentioned has no
atoms at all.  This is not known to cause trouble: a variable with no defined
literal appears in no reason and no conclusion, so \invDomain{} is never
appealed to for it.  But it is a real difference between the framework as
stated here and the framework as implemented, we have not worked the lazy
variant through, and the honest position is that the results below are proved
for \cref{eq:rootbounds} as written.

In the other direction the implementation is \emph{stronger} than \inv{Bound}
asks, and this one costs measurable proof.  \code{fix\_bound} emits a pin for
every out-of-range cut it creates, not one a side, so a variable can accumulate
hundreds of top-of-proof units that \cref{lem:chain} would supply from one.  On
the TSP example this is $2{,}258$ pins over $68$ distinct variable-and-side
pairs, the worst single variable carrying $255$ --- a cost variable with lower
bound $0$ whose cuts run from $\atge{w}{-482}$ up to $\atge{w}{0}$, each pinned
separately.  Weakening the emission to match the invariant is not quite as
simple as pinning the extreme cut, because the cuts are created one at a time
and in no particular order, and the pin has to be a persistent top-of-proof line
rather than one re-derived per use; the workable version pins the cut at the
definition bound once, on the first out-of-range request for the variable, and
nothing after that.
\end{remark}

\begin{theoremT}
\label{thm:maintain}
A correct proof logging CP solver that always logs propagation and conflict
justifications, maintains the literal-layer invariants of \cref{sec:atomprops},
and initially introduces bound literals, will always maintain \invConflict{} and
\invDomain{}.
\end{theoremT}

\begin{proof}
Let $t$ be any point immediately before a backtrack or domain change, with
$B_t$ and $F_t$ as above, and assume \invDomain{} has been maintained at all
relevant earlier points.

Suppose the solver has detected infeasibility of the current domain state.  Then
it will just have logged a conflict justification $R\Rightarrow 0\geq 1$.
\cref{lem:reasons} tells us that reverse unit propagation of the backtracking
justification either reaches contradiction or propagates every literal of $R$.
Either way a contradiction is reached, so \invConflict{} is respected.

Suppose instead a propagator has just made an inference that changed a domain,
and logged $R\Rightarrow\ell\geq 1$.  The reason is valid at the point directly
before the inference, so every literal of $R$, and hence $\ell$, propagates under
$F_t\cup\{\neg B_t\}$.  If $\ell$ is stated over a proper view variable $W$ of
$X$, then by \cref{lem:transport} its image on $X$ propagates too; either way the
domain restriction $\ell$ encodes is incorporated into $\dm_L(X)$ for the
propagated set $L$, since $\dm_L$ is computed over all view variables.  Hence
\invDomain{} continues to be respected.

It remains to argue that there is a starting point at which \invDomain{} holds.
Consider the very first conflict or propagation inference.  The only domain
changes that could have happened are those due to initial decisions, which by
definition propagate on reverse unit propagation of the backtracking
justification, and the initial domain bounds, which we have introduced as units
by \cref{eq:rootbounds}.  Hence \invDomain{} is initially respected.
\end{proof}

\paragraph{Summary of the framework.}
\begin{enumerate}[itemsep=1pt]
  \item Encode the problem using \cref{proc:encoding}, giving every operand
    --- base variable or view --- its own bit vector, and every proper
    view variable its definitional link \cref{eq:link}.
  \item At the start of the proof, derive by RUP $\atge{x}{l}\geq 1$ and
    $\natge{x}{u+1}\geq 1$ for the bounds $l..u$ on each variable, and
    likewise for any out-of-range cut subsequently created on any view variable.
  \item Any time the solver backtracks, derive by RUP a backtracking
    justification $G\Rightarrow 0\geq 1$, where $G$ is the sequence of decisions
    prior to backtracking.
  \item Any time the solver detects infeasibility, derive a conflict
    justification $R\Rightarrow 0\geq 1$.
  \item Any time a propagator makes an inference, derive a propagation
    justification $R\Rightarrow\ell\geq 1$.
  \item Any atomic variables needed for the above must be defined in the proof,
    if not already in the input model, with the corresponding definition
    constraints \cref{eq:defge,eq:defeq,eq:defin}.
  \item Every defined bound literal must be threaded into its view variable's
    order chain (\invChain{}).
  \item Every defined interval literal must be maintained within its
    view variable's partition, with its covering and its containment edges
    (\invPart{}, \invCover{}, \invCont{}).
  \item Every defined atomic literal on a proper view variable, and every one
    on a base that has proper view variables, must be mirrored onto the others
    and joined to them by the linking constraints of
    \cref{eq:gelink,eq:eqlink,eq:inlink} --- for interval literals, triggered at
    the point the literal is \emph{named} (\invView{}, \invName{}).
\end{enumerate}

Steps~1 to~6 are the thesis's framework; steps~7 to~9 are the literal-layer
invariants stated as obligations.  Despite requiring a delicate set-up in terms
of variable encodings and propagation properties, the resulting proof has a
relatively intuitive structure: other than the definition, chain, covering,
containment and linking constraints, it is essentially the solver writing down all
the facts that it learned, at the points that it learned them.

\subsubsection{Optimality Proofs from Solution-Improving Search}
\label{sec:optimality}

There is not a large conceptual difference between complete backtracking search
to establish unsatisfiability and solution-improving search to determine an
optimum.  The only change required from the proof's point of view is that
whenever the solver finds a solution and introduces a new objective bound, it
should use the solution-logging rule to witness a solution, introducing a
constraint requiring objective improvement, define the corresponding bound
literal, and introduce by RUP the constraint setting it to true.

In general, for an objective variable $W$ and attained objective value $o$, a
solution witness introduces
\[
  \BinEnc(W) \geq o+1 \quad\text{for maximisation, or}\quad
  -\BinEnc(W) \geq -o-1 \quad\text{for minimisation,}
\]
followed by a literal definition and the RUP constraint $\atge{w}{o+1}\geq 1$ or
$\natge{w}{o}\geq 1$.  This fits neatly into the framework, as the RUP
constraint derived after a solution is found is simply a propagation
justification with $R=\emptyset$, as if the objective-improving constraint had
always been part of the model.  It has the effect of maintaining \invDomain{} by
ensuring the bound restriction on the objective variable's domain propagates ---
and, by \cref{lem:transport}, on every view variable of it.

\paragraph{A requirement for partial solution witnessing.}
For solution witnessing to be compact and user-friendly, we would like to
witness only a single equality literal for each variable-value pair in the found
assignment to the original CP variables, rather than writing out a complete
assignment to every PB variable in the encoding.  For a partial assignment to be
verifiable as a solution witness, the checker requires that the given assignment
unit propagates to a complete solution of the PB problem.  We guarantee this by
requiring the encoding to have the property that any partial assignment setting
all the bit variables $x_i\in\BinEnc(X)$ for every $X\in\mathcal V$ propagates
to a total assignment or to contradiction.

For the linearisations of the thesis, this is established by observing that any
variable fully reified on a constraint involving only bit variables is
propagated true or false by an assignment to those bit variables, and then
recursively for any variable fully reified on a constraint containing only
variables already known to propagate.  Three additions have to be checked in the
extended encoding, and each extends the same recursion by one step.

\begin{enumerate}[itemsep=1pt]
  \item \textbf{Proper view variables.}  Fixing $\BinEnc(X)$ reduces
    $\ViewLink(W)$, that is $\BinEnc(W) - s_W\BinEnc(X) = c_W$, to a pair of
    constraints of the form $\BinEnc(W)\geq A$ and $\BinEnc(W)\leq A$ for a
    constant $A$.  By Theorem~2.8 of the thesis these unit propagate to a fixed
    consistent assignment of every bit of $\BinEnc(W)$.  So the bits of every
    proper view variable propagate as soon as the base's bits are set.
  \item \textbf{Inequality and equality atoms} on any view variable are fully
    reified over that view variable's bits, or over its inequality atoms, so
    they propagate by the existing recursion once step~1 has fired.
  \item \textbf{Interval atoms} are fully reified, by \cref{eq:defin}, over two
    inequality atoms of the same view variable, so they propagate once those
    do.
\end{enumerate}

None of the atom-level linking constraints of
\cref{eq:gelink,eq:eqlink,eq:inlink} is needed
for this argument, and none of the covering or containment clauses is either;
they are implied clauses over variables that already propagate, and cannot
block a propagation that would otherwise occur.

\subsubsection{Enumeration Proofs from Solution-Excluding Search}
\label{sec:enumeration}

When a solver performs complete solution-excluding search --- adding a nogood
eliminating each solution once found --- the conclusion can be handled in a
manner conceptually similar to proofs of optimality: witness each solution,
relying on unit propagation for completion as above, with a rule that introduces
the negation of the conjunction of literals encoding the solution; then, when
the solver establishes unsatisfiability, conclude that the problem has no
solutions other than those witnessed.

Because the proof system is not implicational --- redundance and dominance can
add or remove solutions --- a formal enumeration conclusion requires a
guarantee that a distinguished set of variables is \emph{preserved}, in the
sense that no rule application may change the projected solution set for them.
The required property of the encoding is then:

\begin{description}[font=\normalfont\bfseries,leftmargin=2.4em]
  \item[P3.] There is a one-to-one correspondence between the solutions to the
    CP problem, and the solutions to the PB problem projected onto a preserved
    set.
\end{description}

The extended encoding makes the choice of preserved set more delicate than it
looks, and we state it explicitly: \textbf{the preserved set is the bits of the
base variables only}.  It does not include the bits of proper view variables,
nor the atomic variables of any view variable, proper or not.  This is the right choice, and
the only workable one, for three reasons.

\begin{itemize}[itemsep=1pt]
  \item It is sufficient.  By \cref{thm:p1p2} the bijection $\phi$ is already
    onto the bits of $\mathcal V$; every other variable in the encoding is
    determined by those, by \cref{lem:composing}.
  \item It is necessary.  Atomic variables are introduced lazily, by redundance,
    during the proof, and the set of them in existence depends on what the
    solver happened to do.  A preserved set that grew as the proof ran would make
    the conclusion depend on the search, which is exactly what a preserved set
    exists to prevent.
  \item It is stable under adding a view variable.  Two runs of the same solver
    that wrap the same variable in different views have different sets of
    proper view variables and therefore different encodings, but the same
    preserved set and the same enumeration conclusion.
\end{itemize}

One consequence deserves care.  A solution witness \emph{is} written over atomic
variables --- one base equality literal per variable --- even though atomic
variables are not preserved.  That is sound only because unit propagation
completes such a witness to the preserved bits, and onward --- through
\cref{sec:optimality}'s step~1 --- to the bits of every proper view variable.
The completion argument is therefore load-bearing for enumeration in a way it is
not for optimality, and a view variable whose bits did not propagate from the
base's would break P3 rather than merely making a witness verbose.

There is a second, quieter interaction.  Some of the units this framework pins
at the top of the proof --- the boundary pins of \inv{Bound} in particular ---
must survive the point at which enumeration restricts the formula, since they
are facts about the encoding rather than about any one solution.  Emitting them
once as persistent top-of-proof lines, rather than re-deriving them at each use,
is what makes that work.

\subsection*{Interlude: what is load-bearing, and how it fails}
\phantomsection\addcontentsline{toc}{subsection}{Interlude: what is load-bearing, and how it fails}\label{sec:loadbearing}

This section is not part of the theory.  It records how the theory fails when
one of its invariants is dropped, because that turns out to be much harder to
observe than one would expect, and because the design has been simplified
incorrectly three times on the strength of tests that stayed green.

\subsubsection*{Two failure modes that look identical}
\phantomsection\label{sec:p1p2}

\begin{description}[font=\normalfont\bfseries,leftmargin=2.4em]
  \item[P1.] A line we emit is not itself accepted by the verifier.  These
    failures are loud, immediate and local: the checker rejects at the line, and
    the error points at the culprit.
  \item[P2.] Every line checks individually, but the clause set does not keep
    unit propagation strong enough for \emph{later} RUP checks --- backtracking
    justifications, other constraints' inferences --- to re-derive the facts they
    depend on.  The eventual rejection lands on an unrelated later line, only
    under composition, and only for particular search shapes.
\end{description}

Everything in \cref{sec:atomprops} beyond the definitions themselves guards
against P2.  \invChain{}, \invCut{}, \invPart{}, \invCover{}, \invCont{},
\invView{} and \invName{} contribute nothing to the validity of any individual
line; they exist so that \cref{thm:complete-family} holds, and
\cref{thm:complete-family} exists so that \cref{lem:reasons} holds, and
\cref{lem:reasons} exists so that a backtracking justification three thousand
lines later is RUP.

The trap follows immediately, and is worth stating as a rule:

\begin{quote}
  \textbf{A RUP check is strictly stronger than the single unit-propagation pass
  that later checks rely on.  Every P2 clause therefore looks deletable under
  P1 testing.}
\end{quote}

Any linking, covering or containment clause can be removed and a local test will
stay green, because the RUP check re-finds the derivation that unit propagation
would have needed the clause for.  ``This clause was never load-bearing in my
tests'' is \emph{always} observable for a P2 clause on small tests.  It is not
evidence.  In particular, a witness stated as a RUP line asserting the fact one
cares about tests nothing at all: only a \emph{backtracking justification},
checked by RUP after the decisions have been asserted, discriminates.

\subsubsection*{The witness suite}

Each clause family has a counterexample that fails quickly if the family is
removed or weakened.  All but one are deterministic tests, checked in as gate-on
and verifier-checked; W4 is instead a standing regression net, the whole
constraint suite run under interval-rejecting branching, and it is random rather
than deterministic.  Any change to the clause set must run them all.

\begin{center}
\small
\begin{tabular}{@{}llp{7.1cm}@{}}
\toprule
& \textbf{Guards} & \textbf{Shape} \\
\midrule
W1 & \cref{def:widthone} & Three chained (dis)equalities on three-value domains;
  branch, then backtrack.  If width-one removals mint interval flags instead of
  using the equality atom, the replay stalls: ``$b$ lost the value $1$'' is
  locked inside a flag with no link to $\ateq{b}{1}$. \\
W2 & \invCover{} & Two variables, one equality and one disequality, with a hole
  run named in a reason.  If that literal cannot be falsified by unit
  propagation, the replay stalls in five steps. \\
W3 & \invCover{}, \invCont{} & A hole run concluded in two pieces and named as
  one interval in a reason.  Containment points the wrong way; only the request
  covering lets the replay through. \\
W4 & \invCont{} & The whole constraint suite run under interval-rejecting
  branching.  Fails within seconds on several counting constraints if the
  containment edges are dropped. \\
W5 & \invCover{}(iii) & A variable whose cells are all excluded must reach
  contradiction at the literal level. \\
W6 & \cref{eq:inlink}, first clause & An interval fact concluded on a base
  variable reaching a reason literal on a view of it. \\
W7 & \cref{eq:inlink}, second clause & The mirror image: concluded on the view,
  needed on the base. \\
W8 & \invName{} & As W6, but the literal named is one the partition machinery
  created as a cell, so it is invisible to a request-driven linking rule. \\
W9 & \cref{eq:inlink} and \invName{} & W8's two-view form, where the fact
  composes $W_1\to X\to W_2$ and the unlinked literal is a conclusion rather
  than a decision. \\
\bottomrule
\end{tabular}
\end{center}

Each of W6 to W9 has been validated by ablation, and every clause and rule is
separately load-bearing:

\begin{center}
\small
\begin{tabular}{@{}lcccc@{}}
\toprule
\textbf{Ablation} & W6 & W7 & W8 & W9 \\
\midrule
none & passes & passes & passes & passes \\
both in-link clauses removed & \textbf{fails} & \textbf{fails} & \textbf{fails} & \textbf{fails} \\
first in-link clause removed & \textbf{fails} & passes & \textbf{fails} & \textbf{fails} \\
second in-link clause removed & passes & \textbf{fails} & passes & \textbf{fails} \\
link on request rather than on naming & passes & passes & \textbf{fails} & \textbf{fails} \\
\bottomrule
\end{tabular}
\end{center}

W2 and W5 are, by contrast, structural: no single-family ablation makes either
fail, because under \invPart{} the vocabulary always matches or grounds out, and
because wipeout is also derivable by the bound-axiom walk of
\cref{rem:rootcovering}.  They remain in the suite as composed end-to-end
regression nets.  Two known gaps: there is no deterministic witness for split
coverings composing across refinements with no bound anchor (the shape
\cref{lem:covering}(b) is proved for), and none for \invCut{} in isolation.

\subsubsection*{The coincidence trap}
\phantomsection\label{sec:coincidence}

A randomised sweep over views is close to no evidence about the in-links, and it
is worth knowing by how much, because the natural reaction to
\cref{rem:bothlinks} is to test it rather than to prove it.

The negative crossing of \cref{eq:inlink} succeeds \emph{with no linking constraint
at all} whenever any of four things happens:

\begin{itemize}[itemsep=1pt]
  \item the interval is the whole definition range, so asserting its negation is
    already inconsistent under unit propagation and the check passes vacuously;
  \item the interval abuts or extends below the lower bound, so that the
    boundary pin of \inv{Bound} --- through the order chain, if the pin is not
    at the cut itself --- makes \cref{eq:din2} a unit, the ge-link carries it,
    and the far side's \emph{forward} reification finishes the job;
  \item the interval abuts or extends above the upper bound, the mirror image;
  \item every cell inside the interval is a singleton --- and one well-placed
    equality atom shatters a width-two cell on \emph{both} sides, so any
    per-value activity anywhere near the interval silently makes the case work.
\end{itemize}

Measured over domain widths $4$ to $16$ with zero to seven random interval
requests per instance: \textbf{80\% of instances cross successfully with no
linking constraint at all} --- abutting above $39\%$, abutting below $23\%$, whole
range $12\%$, fully shattered $5.6\%$ --- and of the $20\%$ that stall, a single
wide cell accounts for $88\%$.  Measured the other way, over $987$ randomly
generated instances that actually reach a conflict, ablating the link constraints
makes only $37$ of them notice.  The sweep is about $96\%$ coincidence.

A discriminating instance therefore needs \emph{all} of: an interval strictly
inside both view variables' definition ranges; at least one interior cell of
width at least two that no per-value machinery touches; and a replay that has to
reach the far side's literal from the near side's negation rather than from a
bound.  W6 to W9 are constructed against exactly this.

\subsubsection*{Corroboration, and what it is worth}

The following are not proofs, and are recorded only for calibration.

\begin{itemize}[itemsep=1pt]
  \item $3{,}204{,}000$ randomised instances with interval inference enabled and
    no verifier failures, once the width-one gap of \cref{def:widthone} was
    closed.  Every failure class encountered before that was reproduced and then
    eliminated by exactly the clauses \cref{sec:atomprops} mandates.
  \item $3{,}000$ randomly generated view proofs with zero \invView{} violations,
    over nine constraint families with holey domains, three-view bases, and
    intervals pushed outside the definition bounds.
  \item Over $8{,}000$ constraint tests under every view wrap and position,
    including mixed wraps of opposite sign in one constraint, verified.
\end{itemize}

Testing doctrine learned at cost, for anyone extending the suite: a capped or
gate-off green run certifies nothing about P2; two-variable instances certify
nothing about composition; and ``the suite passed'' must always be qualified by
which gates and which branching were active.

\subsubsection*{Cost}
\phantomsection\addcontentsline{toc}{subsubsection}{Cost}\label{sec:cost}

Per distinct interval literal requested on a view variable: at most two cell
splits, hence at most four new literals of which any width-one piece is an
equality atom, plus the requested literal; two redundance lines per literal;
binary coverings for each split; one covering for the request, as wide as the
number of prior boundaries inside it; and containment edges to immediate
neighbours.  \textbf{Nothing is proportional to the domain width}, which was the
point.  The known growth mode is a wide request over a heavily fragmented
region, whose covering is as wide as the number of prior boundaries inside it,
bounded by the request count rather than by the domain size.

Per interval literal per proper view variable: the mirrored literal on the
other side, with its own partition maintenance, plus two clauses.  This is the
same constant factor that inequality and equality atoms already pay for a
wrapped variable.  Measured on a constraint over a bare variable of width $w$
and a two-value variable, so that root propagation prunes exactly one wide
interior run:

\begin{center}
\small
\begin{tabular}{@{}rrrrrr@{}}
\toprule
$w$ & bare & view, no interval links & view & negated view & mixed \\
\midrule
$10^3$ & 68 & $71{,}017$ & \textbf{200} & 200 & 200 \\
$10^4$ & 68 & $710{,}017$ & \textbf{200} & 200 & 200 \\
$10^5$ & 68 & \emph{timed out} & \textbf{200} & 200 & 200 \\
$10^6$ & 68 & \emph{timed out} & \textbf{200} & 200 & 200 \\
\bottomrule
\end{tabular}
\end{center}

The third column is the per-value fallback a view had before it owned interval
literals: $17+71w$ lines, unwritable past $10^4$.  The remaining columns are flat.
They are $200$ rather than $68$ because the view carries its own bit vector, its
own atoms and the link constraints --- a constant, not a function of the width.  That
constant is the price of the design; the slope was the point.  The bare column is
the same number on both sides of the change, which is what makes the two
measurements comparable: nothing here touches a variable that no view wraps.

Checking bites before writing does.  The checker takes about $0.012$ seconds on
each flat row, $42$ seconds on the $10^3$ row of the third column, and more than
an hour on the $10^4$ row of that column, at which point the run was stopped.  So
by the width at which the per-value fallback stops being writable, it has already
stopped being checkable.

Two cautions for anyone who repeats this measurement.  The offsets in the mixed
column are not free to choose: the two wraps have to put both operands over the
same interval, because otherwise their ranges are disjoint, the instance is
unsatisfiable at the root, and the row measures a refutation instead of the
interior-hole prune.  An earlier version of this measurement paired $x+10^6$ with
$-y+2\cdot 10^6$, whose ranges meet only at $w=10^6$, so three of the four mixed
rows were the wrong shape and the fourth was the right one; the only symptom was
a three-line step in a column that is otherwise flat.  And the absolute numbers
here are about twice the ones first measured, because every line defining part of
a variable's encoding now also enters the checker's core set, which costs one
further line each.  That scales both columns and leaves flat against linear
untouched.

\appendix
\crefalias{section}{appendix}
\crefname{appendix}{Appendix}{Appendices}
\Crefname{appendix}{Appendix}{Appendices}

\section{A worked example}
\label{sec:worked}

Take $X$ with initial domain $1..20$, wrapped by the single registered view
$V = -X + 30$, so $\varphi_V(k) = 30-k$ and $V$ has definition range $10..29$.
The model contains $\Bnd(X)$, $\Bnd(V)$ and $\ViewLink(V) := V + X = 30$, and
$\BinEnc(X)$ and $\BinEnc(V)$ are separate five-bit vectors.  No interval
literal exists yet, and neither variable has a partition.

\paragraph{First request: $[X\in 5..10]$.}
This may come from any source: a conclusion, a reason element, or a branching
decision.  On $X$:
\begin{itemize}[itemsep=0pt]
  \item the partition is created, with boundaries $\{1,5,11,21\}$ and cells
    $[1,4]$, $[5,10]$, $[11,20]$, each defined as an interval literal with its
    two cuts threaded into the order chain;
  \item the root covering
    $\atin{x}{1}{4}\vee\atin{x}{5}{10}\vee\atin{x}{11}{20}\geq 1$ is emitted;
  \item there is no containment yet, nothing being nested.
\end{itemize}
Because the literal has now been \emph{named} --- by the root covering, if
nothing else --- the mirroring hook fires for each cell in turn, and the images
join $V$'s side.  It is worth following the order, because it is not the order
one would guess.  The first literal the root covering names is $[1,4]$, whose
image $[26,29]$ is what creates $V$'s partition, with boundaries
$\{10,26,30\}$ and cells $[10,25]$, $[26,29]$; $V$'s own root covering is
therefore $\atin{v}{10}{25}\vee\atin{v}{26}{29}\geq 1$, and rendering \emph{that}
names $[10,25]$, whose image is a fresh request $[X\in 5..20]$ back on $X$.  Only
then does $[5,10]$'s own image $[20,25]$ arrive and split $[10,25]$ at $20$.  The
endpoints swap because $\varphi_V$ reverses the order: $30-10=20$, $30-5=25$.
Finally the two clauses
\[
  \natin{x}{5}{10}\vee\atin{v}{20}{25}
  \qquad\text{and}\qquad
  \natin{v}{20}{25}\vee\atin{x}{5}{10}
\]
are emitted.  The recursion terminates because the pair is recorded before the
two \code{need} calls are made.

\paragraph{Second request: $[X\in 7..15]$.}
Both endpoints fall strictly inside existing cells, so both cells split first:
\begin{itemize}[itemsep=0pt]
  \item $[5,10]$ splits at $7$ into $[5,6]$ and $[7,10]$, with the split covering
    $\natin{x}{5}{10}\vee\atin{x}{5}{6}\vee\atin{x}{7}{10}$;
  \item $[11,20]$ splits at $16$ into $[11,15]$ and $[16,20]$, with its split
    covering.
\end{itemize}
Then $[7,15]$ is defined, with the covering
$\natin{x}{7}{15}\vee\atin{x}{7}{10}\vee\atin{x}{11}{15}$ and the two
containment edges $\natin{x}{7}{10}\vee\atin{x}{7}{15}$ and
$\natin{x}{11}{15}\vee\atin{x}{7}{15}$.  The root covering is
\emph{not} touched: to falsify $[5,10]$ later, unit propagation goes through its
own split covering.  The same sequence then happens on $V$ for the image
$[V\in 15..23]$, and the link pair joins the two.

The state now is: boundaries $\{1,5,7,11,16,21\}$ on $X$ and
$\{10,15,20,24,26,30\}$ on $V$; five cells each, in bijection under $\varphi_V$;
and non-cell literals $[X\in 5..10]$, $[X\in 11..20]$, $[X\in 7..15]$ and
$[X\in 5..20]$, with their images.  The last of those is not something any
caller asked for: it is the image of $V$'s first root covering, as above.  This
is the general phenomenon of \cref{sec:open} item~6 --- the two sides hold the
same literals, but arrive at them by different routes and hold them in different
roles.

\paragraph{A negative crossing.}
Suppose the solver now decides $\natin{x}{7}{10}$ --- the \emph{cell} $[7,10]$,
not the request $[7,15]$ --- and, deeper in the subtree, a propagator's reason
names $[V\in 20..23]$, which is its image.  In the replay,
$\natin{x}{7}{10}$ is asserted.  It is a unit on neither of its cuts:
\cref{eq:din2} for it is the ternary clause
$\atin{x}{7}{10}\vee\natge{x}{7}\vee\atge{x}{11}$, and with the first literal
false and neither bound known, nothing propagates.  Its containment edges are no
help either: they run $\natin{x}{7}{10}\vee\atin{x}{5}{10}$ and
$\natin{x}{7}{10}\vee\atin{x}{7}{15}$, both already satisfied by the decision, so
neither propagates anything.  And there is nothing below to descend to: a cell
has no covering, and no defined literal lies strictly inside it.  The only thing
that
carries the fact to $V$ is the first clause of its link pair,
$\natin{v}{20}{23}\vee\atin{x}{7}{10}$: with $\atin{x}{7}{10}$ false it
propagates $\natin{v}{20}{23}$ in one step.  Read as an implication that clause
says $\atin{v}{20}{23}\to\atin{x}{7}{10}$, which is not the direction anything
needs; it is here for its contrapositive.  Delete it and the backtracking
justification fails, while every individual line of the proof still verifies.
This is witness W6 in miniature; W7 is the mirror image --- a fact decided on
$V$ and needed on $X$ --- and it needs the other clause.

\paragraph{Why it has to be a cell.}
Run the same argument on the \emph{request} $[7,15]$ instead and the conclusion
is false: deleting that literal's first link clause leaves the proof verifying.
The reason is that $[7,15]$ is a union of two cells, so the invariants supply a
second route. Containment carries $\natin{x}{7}{15}$ \emph{down} to
$\natin{x}{7}{10}$ and $\natin{x}{11}{15}$; those cells have link pairs of their
own, which carry them to $\natin{v}{20}{23}$ and $\natin{v}{15}{19}$; and $V$'s
request covering
$\natin{v}{15}{23}\vee\atin{v}{15}{19}\vee\atin{v}{20}{23}$ lifts them back up.
So it is the \emph{family} of first clauses that is load-bearing, and no single
member of it, for any literal that is not a cell on both sides.  W6 to W9 are
built on unsplit cells for exactly this reason.

\paragraph{Why a nearby instance proves nothing.}
Take instead the decision $\natin{x}{1}{15}$, whose interval abuts the lower
bound.  Now no link constraint is needed at all.  \cref{eq:din2} for it is
$\atin{x}{1}{15}\vee\natge{x}{1}\vee\atge{x}{16}$; the boundary pin
$\atge{x}{1}$ of \inv{Bound} is a unit, so $\atge{x}{16}$ propagates.  The
ge-link carries that to $\natge{v}{15}$, since $X\geq 16$ is $V\leq 14$, and the
\emph{forward} reification of the image literal $\atin{v}{15}{29}$ then falsifies
it.  A green result on this instance says nothing whatever about the link pair,
and three further shapes have the same property --- see \nsec{sec:coincidence},
where 80\% of small random instances turn out to be one of them.

\section{Residues and open questions}
\label{sec:open}

\begin{enumerate}[itemsep=3pt]

\item \textbf{Variables with no bit vector.}  As noted in
  \nsec{sec:notprovided}, interval literals need cuts to reify against.  The
  gate is whether the variable has a registered bit vector, not what its domain
  looks like: a variable with domain exactly $\{0,1\}$, given a purely direct
  encoding, is registered with a one-element bit vector, and is therefore
  \emph{not} the failing case, even though it has no inequality atoms until
  they are built lazily.  Note that the test is on the two \emph{values}, not
  on the width: a direct-only variable with domain $\{5,6\}$ has no bit vector
  and does throw.  So the failing case is any direct-only variable other than a
  $\{0,1\}$ one.  Consumers guard on the
  same predicate, which resolves a view to its base --- correctly, since a view
  of a bit-vector-less variable would have to mirror onto a variable that cannot
  hold the mirror.

\item \textbf{Initial-domain gaps.}  Holes in an initial domain are currently
  derived at the root of the proof, by root propagation of the constraint that
  posts the domain.  They could instead be stated in the model, as
  $\natin{x}{a}{b}\geq 1$ rows forming part of the domain definition.
  That formulation is equally definitional and slightly more direct, but it
  requires interval literals to exist during model writing, and it changes the
  OPB surface that a verified encoder conforms against.  The decision is open
  and should be taken alongside that work.

\item \textbf{Assertion levels.}  The solver can be asked to \emph{assert} parts
  of the proof rather than derive them, on a scale running
  \textsc{Off} $<$ \textsc{Definitions} $<$ \textsc{Links} $<$
  \textsc{Inferences} $<$ \textsc{Backtracking}.  At \textsc{Links} the atom
  definitions and the linking constraints are asserted, but the covering and
  containment apparatus is still emitted, as ordinary RUP lines: the guards on
  it are all ``above \textsc{Links}''.  Above \textsc{Links} the interval
  machinery is skipped entirely, and no cells exist.  So there is no setting in
  which the links are asserted and the apparatus is dropped, which is what an
  earlier draft of this entry claimed.  The theorems above are stated for the
  fully-derived setting; the asserting levels are not separately exercised by the
  witness suite.

\item \textbf{Unregistered views.}  A view first encountered during proof
  logging, after the model is closed, cannot be given a view variable,
  because \cref{eq:link} is a model constraint.  Such a view is instead rewritten
  onto its base: a bound becomes a bound (with the polarity flip of
  \cref{def:image}), an equality becomes an equality, and an interval becomes
  the image interval with its endpoints swapped when $s=-1$ and re-canonicalised,
  since a width-one image is an equality atom.  This path is sound but it is
  \emph{not} the design of \cref{sec:framework}: nothing about it is in
  $V$-form, and a cutting-planes step that resolved such a literal against a
  constraint body stated over a registered view would strand.  Constants fold to
  the Boolean constants and never reach the literal layer at all.

\item \textbf{Re-entrancy.}  Mirroring an interval literal onto another
  view variable can, in the same call, cause that view variable's partition to
  split, which can mirror back and split this one.  The machinery is therefore
  re-entrant: a boundary is published before the two halves it creates are
  defined, so for that window the partition claims a cell that does not yet
  exist, and definition must be idempotent.  This is one reason \invCover{} is
  stated as a property of the clause set rather than as a rule about when to
  emit what: the emission order is genuinely hard to reason about, whereas the
  structural condition is checkable at any quiescent point.

\item \textbf{Do the two sides really differ?}  We have observed proofs in which
  an interval is a partition cell on one view variable and a request with its
  own covering on its image, and in which the containment DAGs of the two view
  variables are not isomorphic.  We have not characterised when this
  happens.  It does not matter for any result here, and that is deliberate:
  \cref{thm:complete-family} is proved by transporting to one side and then
  arguing there, so it never appeals to a correspondence between the two sides'
  clause sets.  An argument that did appeal to such a correspondence would be
  unsound reasoning even in the runs where the correspondence happens to hold.

\end{enumerate}

\section{Design rationale}
\label{sec:rationale}

The development above says what the encoding is and why it works.  Three of its
choices are not forced, and each has an alternative that a reader will think of;
the arguments are collected here rather than left in the construction, where they
would interrupt it.  Designs that were tried and are \emph{wrong}, as opposed to
merely worse than what we have, are under \cref{sec:refuted}.

\paragraph{A binary encoding of the integer variables.}
The thesis's reason is that $\BinEnc$ is a substitution and not a reformulation:
a linear constraint stays a linear constraint, so P1 and P2 follow from a
bijection between assignments rather than from an argument about what the
constraints mean.  Two further reasons are worth stating, since $\BinEnc$ is
admittedly ``a poor encoding from a solving perspective''.

It is \emph{total}.  Every assignment to the bits denotes an integer, so a
variable's encoding needs no auxiliary constraint to be well defined; $\Bnd(W)$
restricts which integers are allowed, but nothing is needed to make the encoding
mean anything in the first place.  Form preservation alone would not force the
choice, because an order encoding also gives a linear term,
$W = l_W + \sum_{v=l_W+1}^{u_W}\atge{w}{v}$ --- but that identity holds only
where \invChain{} does, and a direct encoding needs an at-most-one for the same
reason.  A bit vector needs nothing.

It is \emph{compact}: $O(\log w)$ PB variables for a variable of width $w$,
against $\Theta(w)$.  Bits are introduced unconditionally, at step~7, for every
member of $\mathcal W$, whereas an atom is introduced only when a constraint or a
propagator asks for one.  Encoding the variables themselves in order or direct
form would therefore put $\Theta(w)$ variables into the formula before the solver
had done anything, and nothing in this design is allowed to be proportional to a
domain's width --- see \nsec{sec:cost}, where that is what the measurement is
for.

The price is propagation, and the design pays it explicitly.  Bit sums propagate
badly: that is what thesis Theorems~2.7 to~2.9 are about, and what
\nsec{sec:notprovided} restates here.  The atom families exist to supply what the
bits cannot express, and the two are incomparable rather than ordered --- a
single bit literal states that $W$ is even, which no atom does, and
$\natge{w}{6}$ states $W\leq 5$, which no bit literal does.

\paragraph{Equality and interval atoms, rather than conjunctions of inequality
atoms.}
Neither needs a Boolean of its own to be definable: $\ateq{w}{v}$ holds exactly
when $\atge{w}{v}$ and $\natge{w}{v+1}$ do, and $\atin{w}{a}{b}$ exactly when
$\atge{w}{a}$ and $\natge{w}{b+1}$ do.  A propagator could write the conjunction
wherever it wants the fact, and the encoding could carry one atom family instead
of three.

The reason it does not is the negation.  Asserting the conjunction is harmless:
both cuts become units either way.  Asserting its negation is not, because
$\neg(\atge{w}{v}\wedge\natge{w}{v+1})$ is the disjunction
$\natge{w}{v}\vee\atge{w}{v+1}$, which is a unit on neither cut.  Removing a
value, or a run of values, is precisely asserting that negation --- so without a
Boolean, the fact a propagator most often wants to state is the one fact unit
propagation cannot use.  With one, the negation is a single literal,
\cref{eq:deq2,eq:din2} are its reverse reifications, and the hole jump that
\cref{thm:wipeout} turns on --- a bound at the lower cut stepping over the whole
hole --- goes through in one step.

Two things follow.  A reason is a set of literals concluding a single literal, so
a propagator that removes a value needs a literal to conclude, and a conjunction
gives it none.  And the same argument at width greater than one is why interval
literals exist: not for strength, since $\atin{w}{a}{b}$ says exactly what its
two cuts say, but because the alternative is $b-a+1$ separate equality atoms, so
that naming a run of holes in a reason would cost the width of the run rather
than one literal.

Stated as a principle: the encoding gives a fact a Boolean of its own exactly
when its \emph{negation} has to be usable as a unit.  That is also why there is
no fourth family for $\leq$ --- $\natge{w}{v}$ already is one.

\paragraph{A view with its own variable, rather than substituted away.}
The alternative is to substitute $sX+c$ for the view wherever it occurs, and
state everything in $X$'s atoms.  It would be sound, and it would save a linear
equality and a bit vector for each view.  We do not do it, for a reason that only
becomes visible in \cref{sec:framework}: it makes the atoms appearing in the
encoding of a constraint depend on \emph{which} wrapper the constraint happened
to be handed, so a generic constraint-logging procedure can no longer treat every
operand identically, and a cutting-planes step that resolves an operand against a
constraint body has to know which view variable that body is stated over.  Giving
the view a variable of its own costs one linear equality and makes every operand
look the same.  The narrower version of the same question --- keep the view's
variable, but state its \emph{intervals} on the base --- is the ``deviewing''
entry of \cref{sec:refuted}, and fails the same way at the level of one atom
kind.

\section{Refuted designs}
\label{sec:refuted}

Each of the following was proposed, looked right, and is wrong.  They are
recorded with the witness that catches each, so that a future simplification
proposal can be tested rather than argued about.  Choices that are merely worse
than the design we have, rather than wrong, are under \cref{sec:rationale}.

\begin{enumerate}[itemsep=3pt]

\item \textbf{``Reify to the two cuts and use the order chain; no covering, no
  containment.''}  A falsification matrix of \emph{derivations} passes --- that
  is P1 only.  Caught by W4 (containment) and W2, W3, W5 (covering).

\item \textbf{``Containment yes, covering unnecessary --- the chain gives those
  deductions.''}  True for RUP-derivability, false for unit propagation.  Caught
  by W2 and W3.

\item \textbf{``An interval removal is one RUP line and the per-value facts
  follow for free.''}  False across a bit-sum equality; needs the explicit
  bridge lemmas under \nsec{sec:notprovided}.

\item \textbf{Coalescing per-value reasons into interval flags after the fact.}
  Mints literals that nothing concludes and nothing covers, hence unfalsifiable
  in a replay; and materialises the per-value equality atoms anyway, through
  reason naming.  Caught by W2 and W3.

\item \textbf{Width-one interval flags.}  Doppelg\"angers of the equality atom;
  see \cref{def:widthone}.  Caught by W1 --- and it bites harder than predicted:
  with width-one flags defined instead of equality atoms, W1 fails \emph{even
  with the coverings and containment intact}, contrary to a pre-implementation
  analysis which held that a flag-to-atom covering would rescue it.

\item \textbf{``The composition failure is specific to one constraint'' /
  ``two variables versus three'' / ``unit propagation cannot thread a bound, so
  a global cross-variable covering $\atge{x}{k}\Leftrightarrow\atge{y}{k}$ is
  required.''}  All artefacts of P2 gaps, diagnosed on populations contaminated
  by the width-one bug.  No cross-variable covering exists in this design and
  none is needed; see \nsec{sec:notprovided}.

\item \textbf{Cover-on-use over the live witness set,} instead of the partition
  invariant.  Sound, and validated, but it drags search state into the encoding
  layer --- a level-aware registry of live conclusions --- and makes the
  completeness argument depend on what was live when.  Rejected for complexity,
  not for soundness.

\item \textbf{Rewriting a view's interval conditions onto its base}
  (``deviewing''), instead of giving the view its own interval literals.  This
  would leave intervals as the one atom kind not stated over the view's own
  variable, so a cutting-planes step resolving an interval literal against
  a constraint body would find nothing to cancel against.  It is the right
  answer only for the unregistered-view path of \cref{sec:open}, which has no
  encoded variable to be consistent with.

\item \textbf{``A registered view needs no interval linking, because an
  interval fact decomposes into: cross one bound, move within one variable,
  cross back.''}  The decomposition is correct for a \emph{positive} interval
  literal, which is a conjunction of two bounds.  It is false for a negated one,
  which is a disjunction.  See \cref{rem:bothlinks}; caught by W6 and W7, and by
  neither a randomised sweep (\nsec{sec:coincidence}) nor a witness stated as a
  RUP line (\nsec{sec:p1p2}).

\end{enumerate}

\paragraph{Change protocol.}  Any proposal to weaken or remove a clause family
from \cref{sec:atomprops} must (a) say which of W1 to W9 it expects to remain
green and why, (b) run the full witness suite with the verifier enabled, and
(c) account for the P1/P2 distinction explicitly.  A locally green test is not
evidence.  Three prior simplifications passed every test their authors thought
to run.

\end{document}
